\documentclass[11pt,a4paper,titlepage]{report}

% --- Packages fondamentaux ---
\usepackage[utf8]{inputenc}
\usepackage[T1]{fontenc}
\usepackage[french,english]{babel}
\usepackage{geometry}
\geometry{top=2.5cm, bottom=2.5cm, left=2.5cm, right=2.5cm}

% --- Mathématiques et symboles ---
\usepackage{amsmath}
\usepackage{amssymb}
\usepackage{amsfonts}
\usepackage{mathtools}

% --- Graphiques et diagrammes TikZ ---
\usepackage{tikz}
\usetikzlibrary{shapes,arrows,arrows.meta,positioning,calc,fit,backgrounds,decorations.pathreplacing,patterns}
\usepackage{pgfplots}
\pgfplotsset{compat=1.18}

% --- Tableaux ---
\usepackage{booktabs}
\usepackage{tabularx}
\usepackage{xltabular}
\usepackage{array}
\usepackage{multirow}
\usepackage{colortbl}

% --- Code source et encadrés ---
\usepackage{listings}
\usepackage{xcolor}

\definecolor{codegreen}{rgb}{0,0.6,0}
\definecolor{codegray}{rgb}{0.5,0.5,0.5}
\definecolor{codepurple}{rgb}{0.58,0,0.82}
\definecolor{backcolour}{rgb}{0.96,0.97,0.98}
\definecolor{darkblue}{rgb}{0.0,0.2,0.6}
\definecolor{alertred}{rgb}{0.8,0.1,0.1}

\lstdefinestyle{cppstyle}{
    backgroundcolor=\color{backcolour},   
    commentstyle=\color{codegreen}\itshape,
    keywordstyle=\color{darkblue}\bfseries,
    numberstyle=\tiny\color{codegray},
    stringstyle=\color{codepurple},
    basicstyle=\ttfamily\footnotesize,
    breakatwhitespace=false,         
    breaklines=true,                 
    captionpos=b,                    
    keepspaces=true,                 
    numbers=left,                    
    numbersep=5pt,                  
    showspaces=false,                
    showstringspaces=false,
    showtabs=false,                  
    tabsize=2,
    frame=single,
    rulecolor=\color{codegray!30}
}
\lstset{style=cppstyle}

% --- Liens et références ---
\usepackage{hyperref}
\hypersetup{
    colorlinks=true,
    linkcolor=darkblue,
    filecolor=magenta,      
    urlcolor=cyan,
    citecolor=darkblue,
    pdftitle={Spécifications Techniques et Documentation d'Ingénierie - Architecture Logicielle PiCar-S FastDrink},
    pdfauthor={Équipe S5-GI FastDrink},
    pdfsubject={Documentation Complète des Services, Bibliothèques, Pilotes et Registres},
    pdfkeywords={PiCar-S, Microservices, Temps Réel, IPC, POSIX SHM, C++17, FSM, HAL, I2C}
}

\usepackage{fancyhdr}
\setlength{\headheight}{28pt}
\pagestyle{fancy}
\fancyhf{}
\fancyhead[L]{\slshape \leftmark}
\fancyhead[R]{\slshape FastDrink S5-GI}
\fancyfoot[C]{\thepage}
\renewcommand{\headrulewidth}{0.4pt}
\renewcommand{\footrulewidth}{0.4pt}

% --- Commandes pratiques ---
\newcommand{\reg}[1]{\texttt{0x#1}}
\newcommand{\file}[1]{\texttt{\detokenize{#1}}}
\newcommand{\code}[1]{\texttt{\detokenize{#1}}}

\begin{document}

% ================= PAGE DE TITRE =================
\begin{titlepage}
    \centering
    \vspace*{1cm}
    
    {\Huge \bfseries FASTDRINK S5-GI}\\[0.5cm]
    {\LARGE \bfseries Manuel d'Ingénierie Système \& Spécification Complète}\\[0.3cm]
    {\Large Architecture Microservices, Registres Matériels, Protocoles IPC et Analyse des Défaillances}\\[1.5cm]
    
    \rule{\linewidth}{0.5mm}\\[0.4cm]
    {\huge \bfseries DOCUMENTATION TECHNIQUE INTÉGRALE}\\[0.2cm]
    {\large Plateforme SunFounder PiCar-S (ARM64 Linux) \& Jumeau Numérique Godot 4}\\[0.4cm]
    \rule{\linewidth}{0.5mm}\\[2.5cm]
    
    \begin{minipage}{0.48\textwidth}
        \begin{flushleft} \large
            \textbf{Projet :} FastDrink S5-GI\\
            \textbf{Système :} PiCar-S / Raspberry Pi\\
            \textbf{Noyau :} Linux RT / Sysfs / POSIX IPC\\
            \textbf{Langage :} C++17 / C++20
        \end{flushleft}
    \end{minipage}
    \begin{minipage}{0.48\textwidth}
        \begin{flushright} \large
            \textbf{Version du document :} 2.4.0\\
            \textbf{Date :} Novembre 2026\\
            \textbf{Statut :} Spécification Formelle\\
            \textbf{Validation :} Banc de Test \& Godot
        \end{flushright}
    \end{minipage}

    \vfill
    {\small Document technique confidentiel élaboré pour la réalisation du véhicule autonome FastDrink S5-GI.\par}
\end{titlepage}

% ================= TABLE DES MATIÈRES =================
\tableofcontents
\newpage

% ================= LISTE DES FIGURES ET TABLEAUX =================
\listoffigures
\listoftables
\newpage

% ================= CHAPITRES =================
\chapter{Architecture Globale \& Interfaces Noyau / Processus}
\label{chap:architecture}

\section{Introduction et Contexte Système}
Le projet FastDrink S5-GI vise la conception d'un véhicule autonome capable de naviguer de manière autonome sur un tracé balisé (ruban noir de largeur 18~mm) tout en maintenant l'intégrité dynamique d'un chargement fragile (une bille métallique reposant au sein d'une coupelle concave peu profonde). 

La contrainte physique fondamentale impose une accélération totale (linéaire et centripète) strictement bornée~:
\begin{equation}
    \|\mathbf{a}\| = \sqrt{a_{\text{tangentiel}}^2 + a_{\text{centripète}}^2} \le a_{\max} = 0{,}735\text{ m/s}^2
\end{equation}
Tout dépassement de ce seuil critique conduit immédiatement au déversement de la bille par franchissement de la bordure ($\Delta h = 1{,}50\text{ mm}$ pour un rayon de cuvette $r = 20\text{ mm}$).

Pour garantir la résilience, le déterminisme temps réel et la testabilité unitaire, le logiciel est articulé autour d'une \textbf{architecture microservices modulaire en C++17}, s'exécutant directement sous Linux (noyau ARM64 sur Raspberry Pi ou x86\_64 en simulation Godot). Chaque brique fonctionnelle est isolée dans un processus indépendant possédant son propre espace d'adressage virtuel, protégeant l'ensemble de la pile contre les fautes de segmentation locales.

\section{Topologie des Processus et Découpage Fonctionnel}
La pile logicielle se décompose en sept services autonomes, orchestrés par un superviseur système faisant office de gestionnaire de cycle de vie~:
\begin{enumerate}
    \item \textbf{\texttt{ServiceManager} (PID 1 / Init)}~: Initialise les points de montage virtuels du noyau (\file{/dev}, \file{/proc}, \file{/sys}, \file{/dev/shm}), configure l'interface réseau Ethernet statique, lance le démon SSH Dropbear et supervise l'ensemble des démons applicatifs avec réarmement automatique.
    \item \textbf{\texttt{FSM} (Finite State Machine / Cœur Décisionnel)}~: Exécute la logique de haut niveau, le régulateur de suivi de trajectoire (PD barycentrique), le profileur cinématique trapézoïdal et les manœuvres d'évitement d'obstacle et de marche arrière guidée.
    \item \textbf{\texttt{Motors} (Contrôle d'Actionnement)}~: Pilote le contrôleur PWM/Servo PCA9685 via le bus I2C et les ponts en H TB6612 via les GPIOs Sysfs. Intègre un filtre de limitation de pente angulaire (slew-rate limiter) pour la direction.
    \item \textbf{\texttt{Sensors} (Acquisition et Traitement du Signal)}~: Échantillonne le module suiveur de ligne à 5 phototransistors TCRT5000 sur bus I2C et le télémètre ultrasonique HC-SR04 via GPIOs à déclenchement impulsionnel. Applique un pipeline de filtrage numérique à double étage (médian glissant + Kalman 1D / EMA).
    \item \textbf{\texttt{PeripheralsService} (Interface Homme-Machine et Événements)}~: Scrute les boutons poussoirs physiques (Start, WiFi) avec anti-rebond matériel/logiciel, module le buzzer piézo-électrique et diffuse les événements d'initialisation sur socket Unix.
    \item \textbf{\texttt{TelemetryService} (Concentrateur de Diagnostic et Diffusion)}~: Agrège l'état global du robot depuis la mémoire partagée, effectue la surveillance dynamique des violations d'accélération, publie un flux TCP unifié et diffuse des trames UDP broadcast sur le réseau local.
    \item \textbf{\texttt{SimBridge} (Passerelle Jumeau Numérique Godot 4)}~: Établit une passerelle réseau TCP bidirectionnelle (port 9000) avec le moteur physique Godot Engine, permettant l'interchangeabilité parfaite entre robot physique et simulateur sans modifier une seule ligne du code de décision.
\end{enumerate}

\section{Mécanismes de Communication Inter-Processus (IPC)}
Pour assurer des transferts à latence sub-microseconde sans surcharge système, trois canaux de communication spécialisés coexistent~:

\subsection{Mémoire Partagée POSIX Globale (\texttt{/dev/shm/picar\_shm})}
Ce segment de mémoire partagée POSIX (\code{shm\_open}, \code{mmap}) stocke la structure \code{RobotTelemetry}. Il permet un échange \emph{zero-copy} bidirectionnel entre la perception (\code{Sensors}), la décision (\code{FSM}) et l'actionnement (\code{Motors}). 
La cohérence des données est protégée par un verrou d'exclusion mutuelle inter-processus POSIX (\code{pthread\_mutex\_t}) initialisé avec l'attribut \code{PTHREAD\_PROCESS\_SHARED}. Un compteur de battements de cœur (\code{heartbeat}) 64 bits est incrémenté périodiquement par la navigation pour servir de chien de garde matériel/logiciel.

\subsection{Bus Mémoire Pub/Sub par Anneau Circulaire (\texttt{MemoryShareComponent})}
En complément de la SHM globale, une infrastructure Pub/Sub sans verrou basée sur des anneaux circulaires mémoire (\file{libs/shared-services/memory}) permet à chaque service d'émettre des trames typées horodatées (\code{SensorsTelemetry}, \code{MotorsComponentTelemetry}, \code{FSMTelemetry}, \code{ServoTelemetry}) à destination du concentrateur de télémétrie.

\subsection{Sockets de Domaine UNIX (\texttt{AF\_UNIX} / \texttt{SOCK\_STREAM})}
Des sockets locaux créés dans \file{/tmp/} assurent le transport fiable des ordres asynchrones, des signaux d'activation et de reconfiguration dynamique en vol~:
\begin{itemize}
    \item \file{/tmp/fsm.sock}~: Contrôle et télécommande du FSM (ex: \code{trigger\_start}, configuration de gains).
    \item \file{/tmp/motors.sock}~: Décalage de trim (\code{STEERING\_OFFSET}), inversion de polarité.
    \item \file{/tmp/sensors.sock}~: Paramètres de filtrage (\code{FILTER\_ALPHA}, seuils ADC).
    \item \file{/tmp/peripherals.sock}~: Diffusion des événements boutons poussoirs et sonneries buzzer.
\end{itemize}

\section{Diagramme d'Architecture Système et Interfaces Noyau}
La Figure~\ref{fig:arch_globale} illustre l'organisation modulaire, le découplage des espaces mémoire, les chemins de communication IPC et l'interaction avec le sous-système matériel via les interfaces du noyau Linux.

\begin{figure}[htbp]
\centering
\begin{tikzpicture}[
    scale=0.92, every node/.style={transform shape},
    box/.style={rectangle, draw=black!70, rounded corners=3pt, thick, align=center, fill=white, inner sep=6pt},
    svcbox/.style={box, fill=blue!8, draw=blue!60, font=\bfseries},
    driverbox/.style={box, fill=orange!12, draw=orange!70, font=\small},
    kernelbox/.style={box, fill=green!10, draw=green!60!black, font=\small},
    ipcbox/.style={box, fill=purple!10, draw=purple!70, dashed, font=\small},
    hwbox/.style={box, fill=red!10, draw=red!70, font=\bfseries},
    arr/.style={-{Stealth[length=2.5mm]}, thick},
    biarr/.style={{Stealth[length=2.5mm]}-{Stealth[length=2.5mm]}, thick},
    dasharr/.style={-{Stealth[length=2.5mm]}, thick, dashed}
]

% Conteneur Utilisateur
\draw[fill=blue!2, draw=blue!30, dashed, rounded corners=6pt] (-7.2, 5.5) rectangle (7.8, -3.2);
\node[anchor=north west, blue!60!black, font=\bfseries\small] at (-7.0, 5.3) {Espace Utilisateur (Userland - Linux Microservices)};

% ServiceManager
\node[svcbox, fill=gray!20, draw=gray!70, minimum width=14.5cm] (sm) at (0.3, 4.5) {
    \textbf{ServiceManager (PID 1 / Init)} \\
    \footnotesize Montage FS (/dev, /proc, /sys, /dev/shm) \textbar{} Superviseur Daemons (waitpid) \textbar{} IP Statique \textbar{} Dropbear SSH
};

% Services applicatifs
\node[svcbox, minimum width=3.2cm] (sensors) at (-5.2, 2.5) {
    \texttt{Sensors}\\
    \scriptsize 50~Hz / 100~kHz\\
    \tiny TCRT5000 + HC-SR04
};

\node[svcbox, minimum width=3.4cm] (fsm) at (-1.3, 2.5) {
    \texttt{FSM}\\
    \scriptsize 50~Hz (Nominal)\\
    \tiny Décision, PD, Profiler
};

\node[svcbox, minimum width=3.2cm] (motors) at (2.6, 2.5) {
    \texttt{Motors}\\
    \scriptsize 100~kHz (Loop)\\
    \tiny TB6612 + PCA9685
};

\node[svcbox, minimum width=2.4cm] (periph) at (6.2, 2.5) {
    \texttt{Peripherals}\\
    \scriptsize Polling 5~ms\\
    \tiny Start/WiFi, Buzzer
};

\node[svcbox, minimum width=3.2cm] (telemetry) at (-3.3, 0.4) {
    \texttt{TelemetryService}\\
    \scriptsize 50~Hz / TCP 91 / UDP 9002\\
    \tiny Agrégation métriques \& Chien de garde
};

\node[svcbox, minimum width=3.2cm] (simbridge) at (4.5, 0.4) {
    \texttt{SimBridge}\\
    \scriptsize 50~Hz / TCP 9000\\
    \tiny Jumeau Godot 4 (JSON / IPC)
};

% Mémoire partagée centrale
\node[ipcbox, minimum width=14.5cm, fill=purple!12, draw=purple!80] (shm) at (0.3, -1.8) {
    \textbf{Segment Mémoire Partagée POSIX : \texttt{/dev/shm/picar\_shm} (\texttt{RobotTelemetry})}\\
    \scriptsize Verrou \texttt{pthread\_mutex\_t} (Process-Shared) \textbar{} Vitesse \& Braquage Cibles \textbar{} Distance filtrée \textbar{} 5x Capteurs optiques \textbar{} Heartbeat Chien de Garde
};

% Liaisons SHM
\draw[biarr, purple!80!black] (sensors.south) -- (-5.2, -1.3);
\draw[biarr, purple!80!black] (fsm.south) -- (-1.3, -1.3);
\draw[biarr, purple!80!black] (motors.south) -- (2.6, -1.3);
\draw[biarr, purple!80!black] (telemetry.south) -- (-3.3, -1.3);
\draw[biarr, purple!80!black] (simbridge.south) -- (4.5, -1.3);

% Sockets Unix
\draw[dasharr, red!70!black] (periph.west) -- node[above, font=\tiny] {/tmp/peripherals.sock} (motors.east);
\draw[dasharr, red!70!black] (fsm.north east) -- ++(0.3, 0.4) -| node[pos=0.25, above, font=\tiny] {/tmp/fsm.sock} (telemetry.north);

% Conteneur Noyau Linux
\draw[fill=green!3, draw=green!40!black, dashed, rounded corners=6pt] (-7.2, -3.7) rectangle (7.8, -5.6);
\node[anchor=north west, green!50!black, font=\bfseries\small] at (-7.0, -3.8) {Noyau Linux (Kernel Space - Drivers \& Sous-Systèmes)};

\node[kernelbox] (i2c_core) at (-4.5, -4.7) {
    \textbf{Bus I2C-1 Driver}\\
    \texttt{/dev/i2c-1} (ioctl, smbus)
};

\node[kernelbox] (gpio_core) at (1.5, -4.7) {
    \textbf{Sous-système GPIO (Sysfs / Libgpiod)}\\
    \texttt{/sys/class/gpio} (Base BCM2835 auto-détectée)
};

\node[kernelbox] (net_core) at (6.0, -4.7) {
    \textbf{Stack Réseau (TCP/IP)}\\
    Sockets AF\_INET / eth0
};

% Conteneur Matériel Physique
\draw[fill=red!2, draw=red!40!black, dashed, rounded corners=6pt] (-7.2, -6.1) rectangle (7.8, -8.0);
\node[anchor=north west, red!60!black, font=\bfseries\small] at (-7.0, -6.2) {Matériel Embarqué PiCar-S (Hardware Layer)};

\node[hwbox] (tcrt) at (-6.0, -7.2) {
    \textbf{TCRT5000 (5 Voies)}\\
    \scriptsize MCU I2C \texttt{0x11} / \texttt{0x0B}\\
    \tiny 10 octets ADC
};

\node[hwbox] (pca) at (-3.0, -7.2) {
    \textbf{PCA9685 PWM}\\
    \scriptsize I2C \texttt{0x40} (50~Hz)\\
    \tiny Servo + Moteurs
};

\node[hwbox] (tb) at (0.2, -7.2) {
    \textbf{Pont TB6612}\\
    \scriptsize DirA: GPIO17\\
    \tiny DirB: GPIO27
};

\node[hwbox] (sonar) at (2.8, -7.2) {
    \textbf{HC-SR04}\\
    \scriptsize Trig: GPIO12\\
    \tiny Echo: GPIO16
};

\node[hwbox] (periph_hw) at (5.8, -7.2) {
    \textbf{Buzzer / Boutons}\\
    \scriptsize Buzzer: GPIO5\\
    \tiny Start: 13, WiFi: 6
};

% Liaisons Noyau vers Matériel
\draw[biarr] (i2c_core.south) -- (tcrt.north);
\draw[biarr] (i2c_core.south) -- (pca.north);
\draw[biarr] (gpio_core.south) -- (tb.north);
\draw[biarr] (gpio_core.south) -- (sonar.north);
\draw[biarr] (gpio_core.south) -- (periph_hw.north);

% Liaisons Utilisateur vers Noyau
\draw[biarr, blue!80!black] (sensors.south west) to[out=230, in=120] (i2c_core.north west);
\draw[biarr, blue!80!black] (sensors.south west) to[out=250, in=140] (gpio_core.north west);
\draw[biarr, blue!80!black] (motors.south east) to[out=300, in=50] (i2c_core.north east);
\draw[biarr, blue!80!black] (motors.south east) to[out=320, in=60] (gpio_core.north);
\draw[biarr, blue!80!black] (periph.south) to[out=270, in=70] (gpio_core.north east);
\draw[biarr, blue!80!black] (telemetry.east) to[out=0, in=130] (net_core.north);
\draw[biarr, blue!80!black] (simbridge.east) to[out=0, in=70] (net_core.north);

\end{tikzpicture}
\caption{Architecture globale de la pile microservices FastDrink et interfaçage avec le noyau Linux.}
\label{fig:arch_globale}
\end{figure}

\section{Couche d'Abstraction Matérielle (HAL)}
Conformément à la spécification technique \textbf{S10} (symétrie algorithmique), les couches décisionnelles du véhicule ne dépendent d'aucune bibliothèque matérielle directe (\code{sysfs}, \code{smbus}, \code{wiringPi}). 
Toutes les interactions matérielles sont encapsulées derrière deux interfaces virtuelles pures~:
\begin{itemize}
    \item \code{IMotorDriver}~: Définit les fonctions d'actionnement \code{set\_speed(float)}, \code{set\_steering(float)}, \code{set\_steering\_offset(float)} et l'arrêt d'urgence \code{emergency\_stop()}.
    \item \code{ISensorDriver}~: Définit les méthodes d'acquisition \code{read\_raw\_ultrasonic\_cm()}, \code{read\_line\_sensors(uint8\_t[5])} et \code{read\_line\_darkness(float[5])}.
\end{itemize}
Cette conception garantit qu'en mode simulation, le service \code{SimBridge} peut injecter les grandeurs physiques virtuelles dans la SHM POSIX sans qu'aucun changement ne soit nécessaire dans l'exécutable \code{FSM}.

\chapter{Pilotes Matériels, Registres I2C \& Communication Noyau}
\label{chap:hardware_drivers}

\section{Sous-Système Matériel et Cartographie des Bus}
Le robot SunFounder PiCar-S repose sur une carte fille de puissance (Robot HATS) connectée au connecteur d'extension GPIO 40 broches du Raspberry Pi. L'interfaçage matériel est pris en charge par le noyau Linux standard à travers deux mécanismes principaux du noyau~:
\begin{enumerate}
    \item \textbf{Le pilote de périphérique caractère I2C (\file{/dev/i2c-1})}~: Permet la communication sérielle synchrone avec le contrôleur PWM PCA9685, le microcontrôleur du suiveur de ligne et le convertisseur analogique-numérique PCF8591.
    \item \textbf{Le sous-système de contrôle GPIO en espace utilisateur (\file{/sys/class/gpio})}~: Offre un contrôle bas niveau déterministe des signaux discrets pour la commande de direction des ponts en H, les signaux impulsionnels du télémètre ultrasonique et les entrées/sorties d'IHM (boutons et buzzer).
\end{enumerate}

Le Tableau~\ref{tab:hardware_map} récapitule l'assignation matérielle complète des périphériques, leurs adresses physiques et les broches BCM associées.

\begin{table}[htbp]
\centering
\footnotesize
\begin{tabularx}{\textwidth}{lccccX}
\toprule
\textbf{Périphérique} & \textbf{Interface} & \textbf{Adresse / Broche} & \textbf{Fréquence} & \textbf{Direction} & \textbf{Rôle Fonctionnel} \\
\midrule
\textbf{PCA9685} & I2C-1 & \reg{40} (\reg{70}) & 50~Hz PWM & OUT & Générateur PWM (Servo direction + Vitesse moteurs) \\
\textbf{Suiveur Ligne} & I2C-1 & \reg{11} (alt \reg{0B}) & 50~Hz / 100~kHz & IN & Lecture ADC 10 bits des 5 phototransistors TCRT5000 \\
\textbf{PCF8591 ADC} & I2C-1 & \reg{48} & Sur demande & IN & Mesure analogique photorésistances \& tension batterie \\
\textbf{TB6612 DirA} & Sysfs GPIO & BCM 17 & Statique & OUT & Sens de rotation Moteur Gauche (1 = AV, 0 = AR) \\
\textbf{TB6612 DirB} & Sysfs GPIO & BCM 27 & Statique & OUT & Sens de rotation Moteur Droit (1 = AV, 0 = AR) \\
\textbf{HC-SR04 Trig} & Sysfs GPIO & BCM 12 & 50~Hz (10~$\mu$s) & OUT & Déclenchement de la salve acoustique ultrasonique \\
\textbf{HC-SR04 Echo} & Sysfs GPIO & BCM 16 & 50~Hz & IN & Mesure de durée du signal d'écho réfléchi \\
\textbf{Buzzer Piézo} & Sysfs GPIO & BCM 5 & 200~Hz polling & OUT & Signalisation sonore (bips de démarrage, wifi, alerte) \\
\textbf{Bouton WiFi} & Sysfs GPIO & BCM 6 & 200~Hz ($\ge 10$~ms) & IN & Bascule du point d'accès / client WiFi (Active LOW) \\
\textbf{Bouton Start} & Sysfs GPIO & BCM 13 & 200~Hz ($\ge 10$~ms) & IN & Déclencheur physique du départ de course (Active LOW) \\
\bottomrule
\end{tabularx}
\caption{Cartographie matérielle des périphériques PiCar-S sous Linux.}
\label{tab:hardware_map}
\end{table}

\section{Détection Dynamique de la Base GPIO du Noyau (\texttt{sysfs})}
Dans les versions récentes du noyau Linux (notamment sur les kernels 5.x et 6.x pour Raspberry Pi 3, 4 et Zero 2W), la numérotation absolue des broches GPIO sous \file{/sys/class/gpio} n'est plus garantie de débuter à 0. Selon le pilote SoC activé (\code{bcm2835-gpio}, \code{pinctrl-bcm2711} ou \code{pinctrl-rp1}), la base peut être décalée à 512 ou toute autre valeur entière.

Pour garantir une portabilité absolue sans recompilation, la classe \code{PiCarHatsSensorDriver} et le service \code{PeripheralsService} implémentent un algorithme de détection automatique de la base GPIO~:
\begin{lstlisting}[language=C++]
static int detect_gpio_base() {
    DIR* dir = opendir("/sys/class/gpio");
    if (!dir) return 0;
    struct dirent* entry;
    int found_base = 0;
    while ((entry = readdir(dir)) != nullptr) {
        if (strncmp(entry->d_name, "gpiochip", 8) == 0) {
            char label_path[256];
            snprintf(label_path, sizeof(label_path), "/sys/class/gpio/%s/label", entry->d_name);
            int lfd = open(label_path, O_RDONLY);
            if (lfd >= 0) {
                char label[64] = {0};
                ssize_t n = read(lfd, label, sizeof(label) - 1);
                close(lfd);
                if (n > 0 && (strstr(label, "bcm2835") || strstr(label, "3f200000.gpio") || strstr(label, "fe200000.gpio"))) {
                    char base_path[256];
                    snprintf(base_path, sizeof(base_path), "/sys/class/gpio/%s/base", entry->d_name);
                    int bfd = open(base_path, O_RDONLY);
                    if (bfd >= 0) {
                        char base_str[32] = {0};
                        ssize_t bn = read(bfd, base_str, sizeof(base_str) - 1);
                        close(bfd);
                        if (bn > 0) {
                            found_base = std::atoi(base_str);
                            break;
                        }
                    }
                }
            }
        }
    }
    closedir(dir);
    return found_base;
}
\end{lstlisting}
Toute broche physique $\text{BCM}_x$ est ainsi mappée sur le numéro de descripteur système réel $\text{actual\_pin} = \text{found\_base} + \text{BCM}_x$.

\section{Contrôleur PWM PCA9685 (Adresse I2C \texttt{0x40})}
Le PCA9685 est un générateur PWM 16 canaux à résolution 12 bits (4096 pas) piloté par bus I2C. Il dispose d'un oscillateur interne typique à 25~MHz et d'un registre de prédivision d'horloge (\code{PRESCALE}).

\subsection{Plan Mémoire des Registres PCA9685}
Le Tableau~\ref{tab:pca9685_regs} décrit les registres exploités par le pilote d'actionnement \code{PiCarHatsMotorDriver} (\file{libs/shared-headers/hal/picar\_i2c\_registers.hpp}).

\begin{table}[htbp]
\centering
\small
\begin{tabularx}{\textwidth}{ccp{3.5cm}X}
\toprule
\textbf{Registre} & \textbf{Adresse} & \textbf{Nom Symbolique} & \textbf{Description \& Rôle} \\
\midrule
\texttt{MODE1} & \reg{00} & \code{picar::i2c::pca9685::MODE1} & Registre de mode 1 (Mise en veille, auto-incrémentation, redémarrage). \\
\texttt{MODE2} & \reg{01} & \code{picar::i2c::pca9685::MODE2} & Registre de mode 2 (Configuration des sorties : Totem-Pole vs Open-Drain). \\
\texttt{SUBADR1-3} & \reg{02} -- \reg{04} & \code{SUBADR1, SUBADR2, SUBADR3} & Sous-adresses I2C secondaires (non utilisées). \\
\texttt{ALLCALLADR} & \reg{05} & \code{ALLCALLADR} & Adresse d'appel général I2C (défaut \reg{70}). \\
\texttt{LED0\_ON\_L} & \reg{06} & \code{LED0\_ON\_L} & Octet de poids faible du tick de passage à l'état HAUT du canal 0. \\
\texttt{LED0\_ON\_H} & \reg{07} & \code{LED0\_ON\_H} & Octet de poids fort (bits 0--3) du tick ON du canal 0. \\
\texttt{LED0\_OFF\_L} & \reg{08} & \code{LED0\_OFF\_L} & Octet de poids faible du tick de passage à l'état BAS du canal 0. \\
\texttt{LED0\_OFF\_H} & \reg{09} & \code{LED0\_OFF\_H} & Octet de poids fort (bits 0--3) du tick OFF du canal 0. \\
\texttt{ALL\_LED\_ON\_L} & \reg{FA} & \code{ALL\_LED\_ON\_L} & Commande globale ON de l'ensemble des 16 canaux. \\
\texttt{ALL\_LED\_OFF\_H} & \reg{FD} & \code{ALL\_LED\_OFF\_H} & Commande globale OFF (Bit 4 = Forçage à 0V pour Failsafe). \\
\texttt{PRESCALE} & \reg{FE} & \code{PRESCALE} & Registre de prédivision d'horloge fixant la fréquence PWM globale. \\
\bottomrule
\end{tabularx}
\caption{Registres principaux du contrôleur PWM I2C PCA9685.}
\label{tab:pca9685_regs}
\end{table}

\subsection{Configuration des Bits de Contrôle (\texttt{MODE1} et \texttt{MODE2})}
\begin{itemize}
    \item \textbf{Bits de \texttt{MODE1}}~:
    \begin{itemize}
        \item \code{RESTART} (\reg{80})~: Autorise le redémarrage logique avec réactivation des sorties.
        \item \code{AUTO\_INC} (\reg{20})~: Active l'auto-incrémentation automatique des adresses de registres lors des écritures séquentielles par bloc I2C.
        \item \code{SLEEP} (\reg{10})~: Place l'oscillateur interne en mode basse consommation. \textbf{Obligatoire} pour autoriser l'écriture dans le registre \code{PRESCALE}.
        \item \code{ALLCALL} (\reg{01})~: Répond aux trames adressées sur \reg{70}.
    \end{itemize}
    \item \textbf{Bits de \texttt{MODE2}}~:
    \begin{itemize}
        \item \code{OUTDRV} (\reg{04})~: Configure les 16 sorties en structure \textbf{Totem-Pole} (Push-Pull, $5\text{V} / 0\text{V}$) plutôt qu'en drain ouvert, assurant des fronts montants francs pour le pilotage des servomoteurs RC et des entrées PWM du TB6612.
    \end{itemize}
\end{itemize}

\subsection{Formules Mathématiques de Calcul des Registres}
\subsubsection{1. Calcul du Prescaler d'Horloge}
Pour obtenir une fréquence de commande de servomoteurs RC standard de $f_{\text{PWM}} = 50\text{ Hz}$ à partir de l'horloge $f_{\text{osc}} = 25\text{ MHz}$ et d'une résolution de 4096 points~:
\begin{equation}
    \text{prescale} = \text{round}\left(\frac{f_{\text{osc}}}{4096 \times f_{\text{PWM}}}\right) - 1 = \text{round}\left(\frac{25 \times 10^6}{4096 \times 50}\right) - 1 = 122 - 1 = 121 \quad (\texttt{0x79})
\end{equation}

\subsubsection{2. Conversion d'Angle de Braquage en Ticks PWM}
Le servomoteur de direction est alimenté sur une période $T = 20\text{ ms}$ (soit 4096 ticks).
\begin{itemize}
    \item Neutre ($0^\circ$) $\to \tau = 1{,}50\text{ ms} \implies \text{ticks} = \frac{1{,}50}{20} \times 4096 \approx 307\text{ ticks}$.
    \item Braquage Gauche maximal ($-45^\circ$) $\to \tau = 1{,}00\text{ ms} \implies \text{ticks} \approx 205\text{ ticks}$.
    \item Braquage Droit maximal ($+45^\circ$) $\to \tau = 2{,}00\text{ ms} \implies \text{ticks} \approx 410\text{ ticks}$.
\end{itemize}
L'équation d'interpolation affine bornée implémentée dans \code{steering\_angle\_to\_pwm} est~:
\begin{equation}
    \text{ticks}(\theta) = \text{clamp}\left(307{,}0 + \frac{\text{clamp}(\theta, -45^\circ, +45^\circ)}{45{,}0} \times 102{,}0, \; 205, \; 410\right)
\end{equation}

\subsubsection{3. Conversion du Pourcentage de Vitesse en Ticks PWM}
La puissance moteur $[0{,}0\%, 100{,}0\%]$ est transformée en rapport cyclique de 0 à 4095~:
\begin{equation}
    \text{ticks}(P) = \text{clamp}(|P|, 0{,}0, 100{,}0) \times 40{,}95
\end{equation}

\subsubsection{4. Coupure d'Urgence Immédiate (Failsafe Matériel)}
Pour désexciter instantanément tous les canaux PWM sans effectuer 16 transactions individuelles, le pilote injecte la valeur \reg{10} dans le registre \code{ALL\_LED\_OFF\_H} (\reg{FD})~:
\begin{lstlisting}[language=C++]
void PiCarHatsMotorDriver::emergency_stop() {
    m_i2c.write_byte_data(picar::i2c::pca9685::ALL_LED_OFF_H, 0x10); // Bit 4 = Full OFF
}
\end{lstlisting}

\section{Pont en H Double TB6612FNG (Propulsion Arrière)}
Le circuit intégré Toshiba TB6612FNG commande les deux moteurs à courant continu du train arrière. La vitesse de rotation est modulée par les canaux PWM 4 (Moteur Gauche) et 5 (Moteur Droit) du PCA9685. Le sens de rotation est contrôlé par deux broches logiques directes du Raspberry Pi~:
\begin{itemize}
    \item \textbf{Moteur Gauche (\code{DirA})}~: BCM 17 (\file{/sys/class/gpio/gpio17/value}).
    \item \textbf{Moteur Droit (\code{DirB})}~: BCM 27 (\file{/sys/class/gpio/gpio27/value}).
\end{itemize}

Le Tableau~\ref{tab:tb6612_truth} présente la table de vérité opérationnelle appliquée par le pilote.

\begin{table}[htbp]
\centering
\begin{tabular}{ccccc}
\toprule
\textbf{Consigne Vitesse} & \textbf{DirA (BCM 17)} & \textbf{DirB (BCM 27)} & \textbf{PWM (CH 4 \& 5)} & \textbf{Comportement Véhicule} \\
\midrule
$v > 0$ & 1 (HIGH) & 1 (HIGH) & Ticks proportionnels & Marche Avant Nominale \\
$v < 0$ & 0 (LOW) & 0 (LOW) & Ticks proportionnels & Marche Arrière Guidée \\
$v = 0$ & Inchangé & Inchangé & 0 tick & Roue libre / Arrêt doux \\
Arrêt Urgence & 0 (LOW) & 0 (LOW) & 0 tick (Forcé \reg{FD}) & Coupure Totale d'Alimentation \\
\bottomrule
\end{tabular}
\caption{Table de vérité du pont en H double TB6612FNG.}
\label{tab:tb6612_truth}
\end{table}

\section{Module Suiveur de Ligne 5 Voies TCRT5000 (Adresses \texttt{0x11} / \texttt{0x0B})}
Le capteur de détection de trajectoire comporte cinq phototransistors infrarouges TCRT5000 alignés transversalement sur une barrette frontale, espacés d'environ $19\text{ mm}$. Un microcontrôleur embarqué échantillonne les tensions analogiques de réflexion et met à disposition un registre I2C de lecture continue.

\subsection{Protocole de Lecture et Décodage des Trames}
La lecture est initiée par une lecture en bloc de 10 octets consécutifs à partir du registre \code{REG\_START} (\reg{00})~:
\begin{equation}
    \text{Trame I2C} = \begin{bmatrix} H_0 & L_0 & H_1 & L_1 & H_2 & L_2 & H_3 & L_3 & H_4 & L_4 \end{bmatrix}
\end{equation}
Chaque capteur $i \in \{0, 1, 2, 3, 4\}$ dispose d'une mesure analogique sur 10 bits ($0 \le \text{ADC}_i \le 1023$) recomposée selon~:
\begin{equation}
    \text{ADC}_i = (H_i \ll 8) \;|\; L_i
\end{equation}

\subsection{Transformation en Noirceur Continue et Calibration}
Sur un revêtement blanc réfléchissant, les photodiodes reçoivent un flux lumineux élevé, générant une tension et une valeur ADC élevées ($\text{ADC}_{\text{blanc}} \approx 80 \text{ à } 140$). Sur le ruban isolant noir absorbant, le signal chute drastiquement ($\text{ADC}_{\text{noir}} \approx 15 \text{ à } 40$).

La fonction \code{adc\_to\_darkness} projette de manière continue la lecture brute sur l'intervalle $[0{,}0, 1{,}0]$ (où $1{,}0$ désigne le noir absolu)~:
\begin{equation}
    D_i = \text{clamp}\left(\frac{\text{ADC}_{\text{blanc}, i} - \text{ADC}_i}{\text{ADC}_{\text{blanc}, i} - \text{ADC}_{\text{noir}, i}}, \; 0{,}0, \; 1{,}0\right)
\end{equation}
Les seuils de calibration peuvent être personnalisés individuellement par capteur via le fichier persistant \file{/etc/picar\_line\_calibration} ou les variables d'environnement \code{PICAR\_LINE\_BLACK\_ADC} et \code{PICAR\_LINE\_WHITE\_ADC}.

\subsection{Tolérance aux Pannes I2C}
En cas de déconnexion transitoire ou de parasites sur le bus I2C, le pilote tolère jusqu'à 2 lectures erronées consécutives en conservant la dernière valeur acquise valide. Si le compteur dépasse 3 erreurs consécutives ($\approx 60\text{ ms}$), le pilote applique un repli sécuritaire (\emph{failsafe}) en forçant toutes les valeurs à zéro ($D_i = 0$), déclenchant ainsi immédiatement la gestion de perte de ligne par le FSM.

\section{Télémètre Ultrasonique HC-SR04 (GPIOs Sysfs 12 \& 16)}
Le capteur HC-SR04 mesure la distance par temps de vol d'une onde acoustique à 40~kHz.
\begin{enumerate}
    \item \textbf{Impulsion Trigger (BCM 12)}~: Le pilote génère une impulsion haute calibrée d'au moins $10\text{ }\mu\text{s}$ sur la broche \code{actual\_trig}.
    \item \textbf{Mesure Echo (BCM 16)}~: La broche \code{actual\_echo} passe à l'état HAUT pendant la durée de propagation aller-retour.
\end{enumerate}

\subsection{Equation Acoustique et Gestion des Délais}
La vitesse du son dans l'air à $20^\circ\text{C}$ étant $c \approx 343\text{ m/s} = 0{,}0343\text{ cm/}\mu\text{s}$, la distance en centimètres est donnée par~:
\begin{equation}
    d_{\text{cm}} = \frac{\Delta t_{\mu\text{s}} \times c}{2} = \frac{\Delta t_{\mu\text{s}}}{58{,}3} \approx \frac{\Delta t_{\mu\text{s}}}{58}
\end{equation}

\subsection{Protection Contre les Opérations Bloquantes}
Pour préserver les contraintes temps réel strictes sans bloquer le thread d'acquisition~:
\begin{itemize}
    \item \textbf{Chien de garde de front montant}~: Si le signal \code{Echo} ne monte pas dans les 5~ms suivant l'impulsion Trigger, la mesure est avortée avec le code $-1{,}0\text{ cm}$ (obstacle absent ou capteur déconnecté).
    \item \textbf{Chien de garde de front descendant}~: La boucle d'attente s'interrompt automatiquement après 25~ms, ce qui correspond à la portée utile maximale du robot ($\approx 430\text{ cm}$).
\end{itemize}

\section{Diagramme Synthétique de l'Interfaçage Matériel}
La Figure~\ref{fig:hw_schematic} expose le câblage direct entre le microprocesseur BCM2835 et les périphériques de puissance et de détection.

\begin{figure}[htbp]
\centering
\begin{tikzpicture}[
    scale=0.78, every node/.style={transform shape},
    box/.style={rectangle, draw=black!70, thick, rounded corners=3pt, align=center, fill=white, inner sep=6pt},
    rpi/.style={box, fill=red!15, draw=red!70!black, font=\bfseries},
    chip/.style={box, fill=blue!10, draw=blue!70, font=\bfseries},
    sensor/.style={box, fill=green!12, draw=green!60!black, font=\bfseries},
    actuator/.style={box, fill=orange!15, draw=orange!70!black, font=\bfseries},
    bus/.style={-{Stealth[length=2.5mm]}, thick}
]

% RPi Box
\node[rpi, minimum width=4.5cm, minimum height=6.5cm, text width=4.2cm] (rpi) at (-4, 0) {
    \centering\textbf{Raspberry Pi}\\
    \textbf{SoC BCM2835}\\
    \footnotesize (Linux Kernel 5.x/6.x)\\[0.4cm]
    \raggedright\scriptsize
    $\bullet$ I2C-1 SDA (GPIO 2 / Pin 3)\\
    $\bullet$ I2C-1 SCL (GPIO 3 / Pin 5)\\
    $\bullet$ TB6612 DirA (GPIO 17 / Pin 11)\\
    $\bullet$ TB6612 DirB (GPIO 27 / Pin 13)\\
    $\bullet$ Sonar Trig (GPIO 12 / Pin 32)\\
    $\bullet$ Sonar Echo (GPIO 16 / Pin 36)\\
    $\bullet$ Buzzer Out (GPIO 5 / Pin 29)\\
    $\bullet$ Bouton WiFi (GPIO 6 / Pin 31)\\
    $\bullet$ Bouton Start (GPIO 13 / Pin 33)
};

% I2C Bus line
\draw[line width=1.5pt, blue!70] (-1.5, 2.2) -- (4.5, 2.2) node[midway, above, font=\footnotesize\bfseries] {Bus I2C-1 (3.3V)};
\draw[bus, blue!70] (-1.7, 2.2) -- (-1.5, 2.2);

% I2C Chips
\node[chip] (pca) at (1.5, 3.2) {
    \textbf{PCA9685 PWM (0x40)}\\
    \scriptsize 16 Canaux 12 bits \textbar{} 50~Hz
};
\draw[bus, blue!70] (1.5, 2.2) -- (pca.south);

\node[sensor] (tcrt) at (4.5, 3.2) {
    \textbf{TCRT5000 MCU (0x11)}\\
    \scriptsize 5 Voies Suiveur \textbar{} 10 octets ADC
};
\draw[bus, blue!70] (4.5, 2.2) -- (tcrt.south);

% Actuators from PCA
\node[actuator] (servo) at (1.5, 4.8) {
    \textbf{Servo Direction}\\
    \scriptsize Canal 0 PWM (1.0--2.0 ms)
};
\draw[bus, orange!80!black] (pca.north) -- (servo.south);

\node[chip] (tb) at (1.5, 0.5) {
    \textbf{TB6612FNG}\\
    \scriptsize Pont en H Double
};
\draw[bus, orange!80!black] (pca.south) -- node[right, font=\tiny] {CH4/5 PWM} (tb.north);

% Direct GPIO lines
\draw[bus, green!60!black] (-1.7, 0.6) -- node[above, font=\tiny] {DirA (GPIO17)} (tb.west);
\draw[bus, green!60!black] (-1.7, 0.2) -- node[below, font=\tiny] {DirB (GPIO27)} (tb.south west);

\node[actuator] (dc_motors) at (4.8, 0.5) {
    \textbf{Moteurs DC Arrière}\\
    \scriptsize Roues Motrices Gauche \& Droite
};
\draw[bus, orange!80!black] (tb.east) -- (dc_motors.west);

% HC-SR04
\node[sensor] (sonar) at (1.5, -1.8) {
    \textbf{Sonar HC-SR04}\\
    \scriptsize Télémètre Ultrasonique
};
\draw[bus, green!60!black] (-1.7, -1.6) -- node[above, font=\tiny] {Trig (GPIO12)} (sonar.north west);
\draw[bus, green!60!black] (sonar.south west) -- node[below, font=\tiny] {Echo (GPIO16)} (-1.7, -2.0);

% Peripherals
\node[box, fill=purple!10, draw=purple!70] (periph_box) at (1.5, -3.8) {
    \textbf{IHM \& Signalisation}\\
    \scriptsize Buzzer (GPIO5) \textbar{} Start (GPIO13) \textbar{} WiFi (GPIO6)
};
\draw[bus, purple!80!black] (-1.7, -3.8) -- (periph_box.west);

\end{tikzpicture}
\caption{Schéma fonctionnel des connexions physiques et adressage des périphériques.}
\label{fig:hw_schematic}
\end{figure}

\chapter{Spécification Détaillée des Bibliothèques Partagées}
\label{chap:shared_libraries}

\section{Organisation et Statut des Bibliothèques}
L'écosystème logiciel FastDrink repose sur deux ensembles de bibliothèques distincts~:
\begin{enumerate}
    \item \textbf{\texttt{libs/shared-headers/}} (\textbf{GELÉ / INTÉGRITÉ STRICTE})~: Contient l'ensemble des définitions de types fondamentaux, formats de trames binaires, drapeaux d'état, interfaces HAL et algorithmes de régulation mathématique. \textbf{Cette bibliothèque est strictement en lecture seule} afin de garantir la pérennité de l'interface IPC binaire entre tous les processus.
    \item \textbf{\texttt{libs/shared-services/}}~: Fournit l'implémentation opérationnelle des abstractions système partagées (serveur/abonnés de mémoire partagée cyclique, serveurs de sockets UNIX et gestionnaire de fichiers de configuration).
\end{enumerate}

\section{Protocole Binaire et En-tête de Trames (\texttt{header/})}
Toutes les communications sérialisées (sur socket UNIX ou TCP) partagent un en-tête binaire commun de 8 octets à alignement strict sur 32 bits, défini dans \file{header/header.hpp}.

\subsection{Structure \texttt{PacketHeader}}
\begin{lstlisting}[language=C++]
struct PacketHeader {
    HeadersType type;        // 1 octet  : Command (0x01), Telemetry (0x02), Response (0x03)
    uint8_t     headerFlag;  // 1 octet  : Masque de bits (REPLY_REQUESTED, CRC_ENABLED)
    uint16_t    seqNumber;   // 2 octets : Numero de sequence sequentiel (detection de perte)
    uint32_t    payloadLength;// 4 octets: Longueur du payload (DOIT etre multiple de 4)
};
\end{lstlisting}

\subsection{Drapeaux d'En-tête (\texttt{shared::headers::flags})}
\begin{itemize}
    \item \code{REPLY\_REQUESTED} (\reg{01})~: L'émetteur attend une confirmation immédiate (\code{Ack} ou \code{Nack}).
    \item \code{FIRE\_AND\_FORGET} (\reg{02})~: Envoi asynchrone sans attente de réponse.
    \item \code{CRC\_ENABLED} (\reg{04})~: La trame se termine par un mot de contrôle d'intégrité CRC-32 de 4 octets.
\end{itemize}

\subsection{Algorithme de Sérialisation et Contrôle d'Intégrité}
La fonction \code{serializePacket} applique un contrôle strict sur la longueur du payload~:
\begin{equation}
    \text{payloadLength} \pmod 4 = 0
\end{equation}
Si le drapeau \code{CRC\_ENABLED} est actif, la trame sérialisée est prolongée d'un mot de 32 bits calculé selon la norme IEEE 802.3 (\file{crc/crc_global.hpp}) avec le polynôme inverse standard~:
\begin{equation}
    P(x) = \texttt{0xEDB88320}, \quad \text{Valeur initiale} = \texttt{0xFFFFFFFF}
\end{equation}

\section{Segment Mémoire Partagée POSIX (\texttt{ipc/shared\_memory.hpp})}
La structure \code{RobotTelemetry} constitue la mémoire centrale d'échange à latence sub-microseconde entre \code{Sensors}, \code{FSM} et \code{Motors}.

\begin{table}[htbp]
\centering
\small
\begin{tabularx}{\textwidth}{lp{2.5cm}cX}
\toprule
\textbf{Champ Membre} & \textbf{Type} & \textbf{Taille} & \textbf{Description \& Rôle Fonctionnel} \\
\midrule
\code{filtered\_dist} & \texttt{float} & 4 octets & Distance acoustique filtrée (en cm, Médian + Kalman/EMA). \\
\code{speed} & \texttt{float} & 4 octets & Vitesse estimée actuelle du véhicule (en m/s ou cm/s). \\
\code{line\_error} & \texttt{float} & 4 octets & Écart barycentrique latéral de la ligne $[-1{,}0, +1{,}0]$. \\
\code{line\_sensors[5]} & \texttt{uint8\_t[5]} & 5 octets & État logique binaire binarisé des 5 TCRT5000 (1 = ligne noire). \\
\code{line\_darkness[5]} & \texttt{float[5]} & 20 octets & Mesure de noirceur continue normalisée $[0{,}0 = \text{blanc}, 1{,}0 = \text{noir}]$. \\
\code{target\_speed} & \texttt{float} & 4 octets & Consigne de vitesse issue du MotionProfiler (écrite par FSM). \\
\code{target\_steering} & \texttt{float} & 4 octets & Consigne de braquage en degrés (écrite par le PD FSM). \\
\code{steering\_offset} & \texttt{float} & 4 octets & Trim logiciel de correction angulaire mécanique (en degrés). \\
\code{state} & \texttt{uint32\_t} & 4 octets & État machine FSM en cours (cast de \code{MachineStates}). \\
\code{emergency\_stop} & \texttt{bool} & 1 octet & Drapeau d'arrêt d'urgence absolu (Failsafe logiciel immédiat). \\
\code{heartbeat} & \texttt{uint64\_t} & 8 octets & Compteur incrémenté à 50~Hz servant de chien de garde temporel. \\
\code{mtx} & \texttt{pthread\_mutex\_t} & 40 octets & Mutex POSIX inter-processus (\code{PTHREAD\_PROCESS\_SHARED}). \\
\code{cfg\_max\_acceleration} & \texttt{float} & 4 octets & Accélération maximale bornée ($a \le 0{,}735\text{ m/s}^2$). \\
\code{cfg\_cruise\_speed} & \texttt{float} & 4 octets & Vitesse de croisière nominale en ligne droite (m/s). \\
\code{cfg\_min\_curve\_speed} & \texttt{float} & 4 octets & Vitesse plancher en courbe serrée $R = 120\text{ mm}$ (m/s). \\
\code{cfg\_line\_kp} & \texttt{float} & 4 octets & Gain proportionnel du régulateur de trajectoire. \\
\code{cfg\_line\_kd} & \texttt{float} & 4 octets & Gain dérivé du régulateur de trajectoire. \\
\code{cfg\_curve\_safety\_factor} & \texttt{float} & 4 octets & Coefficient de sécurité dynamique virage ($0{,}85 \text{ à } 0{,}93$). \\
\code{cfg\_predictive\_brake\_gain}& \texttt{float} & 4 octets & Gain feedforward sur la dérivée d'erreur pour anticipation. \\
\code{cfg\_version} & \texttt{uint32\_t} & 4 octets & Compteur de version incrémenté à chaque modification en vol. \\
\bottomrule
\end{tabularx}
\caption{Organisation détaillée de la structure mémoire partagée \texttt{RobotTelemetry}.}
\label{tab:robot_telemetry}
\end{table}

\section{Ensemble de Commandes \& Réponses (\texttt{commands/} \& \texttt{responses/})}
Les commandes asynchrones transmises au format \code{PacketHeader} embarquent un sous-en-tête de commande \code{CommandsHeader}~:
\begin{lstlisting}[language=C++]
enum class ServiceID : uint8_t { TELEMETRY = 0x00, SENSORS = 0x01, MOTORS = 0x02, FSM = 0x03 };
enum class CommandID : uint8_t { Ping = 0x01, SetActive = 0x02, SetControlFlags = 0x03, SetConfigKey = 0x04 };

struct CommandsHeader {
    ServiceID serviceID;     // Service cible
    CommandID commandID;     // Type d'action
    uint16_t  payloadLen;    // Longueur des arguments
};
\end{lstlisting}
Pour la commande \code{SetConfigKey}, le payload est constitué de la structure \code{SetConfigPayload}~:
\begin{lstlisting}[language=C++]
struct SetConfigPayload {
    char key[64];            // Cle alphanumerique (ex: "LINE_KP", "MAX_ACCELERATION")
    float value;             // Valeur flottante a appliquer
};
\end{lstlisting}
En retour, le service interrogé émet une réponse \code{ResponseID::Ack} (\reg{02}) ou \code{ResponseID::Nack} (\reg{03}) accompagnée d'un code de rejet \code{NackResponse}~:
\begin{lstlisting}[language=C++]
struct NackResponse {
    uint16_t explanation_len;
    uint16_t nackCode;       // 1 = Erreur inconnue, 2 = Cle invalide, 3 = Valeur hors bornes
};
\end{lstlisting}

\section{Algorithmes de Contrôle et Cinématique (\texttt{control/})}

\subsection{Profileur de Vitesse et Limiteur d'Accélération (\texttt{MotionProfiler})}
Le module \code{MotionProfiler} implémente un générateur de rampe cinématique trapézoïdale garantissant le respect absolu de la contrainte FastDrink. 

\subsubsection{1. Évolution Temporelle Discrète}
À chaque pas temporel $\Delta t$, la vitesse $v_k$ est ajustée selon la consigne désirée $v_{\text{cible}}$ sans jamais dépasser le quantum $\Delta v_{\max} = a_{\max} \cdot \Delta t$~:
\begin{equation}
    v_k = \begin{cases}
        \min(v_{\text{cible}}, \; v_{k-1} + a_{\max} \cdot \Delta t) & \text{si } v_{\text{cible}} > v_{k-1} \quad \text{(Phase d'accélération)} \\
        \max(v_{\text{cible}}, \; v_{k-1} - a_{\max} \cdot \Delta t) & \text{si } v_{\text{cible}} < v_{k-1} \quad \text{(Phase de décélération)}
    \end{cases}
\end{equation}

\subsubsection{2. Vitesse Limite en Virage (Modèle Cinématique d'Ackermann)}
Soit $L = 0{,}14\text{ m}$ l'empattement du véhicule et $\delta$ l'angle de braquage des roues directrices avant. Le rayon de courbure instantané $R$ du centre de l'essieu arrière est~:
\begin{equation}
    R = \frac{L}{\tan(|\delta|)}
\end{equation}
L'accélération centrifuge s'écrivant $a_c = \frac{v^2}{R} \le a_{\max}$, la vitesse critique maximale admissible en virage est~:
\begin{equation}
    v_{\text{virage}}(\delta) = \eta_{\text{safety}} \cdot \sqrt{a_{\max} \cdot \frac{L}{\tan(|\delta|)}}
\end{equation}
où $\eta_{\text{safety}} \in [0{,}85, 0{,}95]$ représente le facteur de sécurité paramétrable (\code{CURVE\_SAFETY\_FACTOR}).

\subsubsection{3. Freinage Prédictif par Anticipation}
Pour éviter les à-coups violents à l'entrée soudaine d'un virage, un terme de freinage prédictif feedforward atténue la consigne de vitesse proportionnellement à la vitesse de variation de l'erreur latérale~:
\begin{equation}
    v_{\text{adapt}} = v_{\text{virage}} - K_{\text{brake}} \cdot \left|\frac{de}{dt}\right|
\end{equation}

\subsection{Asservissement Suivi de Trajectoire (\texttt{LineFollowerController})}
L'orientation du train avant est régulée par un contrôleur proportionnel-dérivé-intégral (PID) barycentrique optimisé.

\subsubsection{1. Calcul d'Écart Barycentrique Analogique}
Chaque capteur $i \in \{0, 1, 2, 3, 4\}$ est affecté d'un coefficient spatial $W = \{-2{,}0, -1{,}0, 0{,}0, 1{,}0, 2{,}0\}$. L'erreur latérale normalisée $e \in [-1{,}0, +1{,}0]$ est calculée par pondération de la noirceur continue $D_i$~:
\begin{equation}
    e = \frac{\sum_{i=0}^{4} W_i \cdot D_i}{2{,}0 \cdot \sum_{i=0}^{4} D_i}
\end{equation}
Si la somme totale $\sum D_i < 0{,}25$, le robot est considéré en perte totale de ruban.

\subsubsection{2. Équation de Contrôle de l'Angle de Braquage}
L'angle de braquage commandé au servomoteur (en degrés) est~:
\begin{equation}
    \theta(t) = K_p \cdot e(t) + K_d \cdot \frac{de(t)}{dt} + K_i \int e(\tau) \, d\tau
\end{equation}
Le terme intégral n'est actif qu'au voisinage immédiat du centre ($|e| < 0{,}6$) afin de compenser les dérives mécaniques de parallélisme sans provoquer d'effet de saturation (\emph{wind-up}) en virage prononcé.

\section{Pipeline de Traitement du Signal (\texttt{filters/})}
Les lectures du capteur ultrasonique HC-SR04 subissent un traitement numérique à deux étages en cascade dans la classe \code{UltrasonicPipeline}~:

\subsection{Étage 1 : Filtre Médian Glissant à 5 Échantillons (\texttt{MedianFilter})}
Le filtre médian élimine strictement les bruits impulsifs, réflexions acoustiques multiples et lectures fantômes à 0~cm~:
\begin{equation}
    y_1[k] = \text{médiane}\Big(x[k], x[k-1], x[k-2], x[k-3], x[k-4]\Big)
\end{equation}

\subsection{Étage 2 : Filtre de Kalman 1D Compensé par la Vitesse}
Le signal débruité alimente un filtre de Kalman 1D dont l'équation de prédiction intègre la vitesse réelle du véhicule~:
\begin{align}
    \hat{d}_{k|k-1} &= \hat{d}_{k-1|k-1} - v \cdot \Delta t \\
    P_{k|k-1} &= P_{k-1|k-1} + Q \\
    K_k &= \frac{P_{k|k-1}}{P_{k|k-1} + R} \\
    \hat{d}_{k|k} &= \hat{d}_{k|k-1} + K_k \cdot \Big(y_1[k] - \hat{d}_{k|k-1}\Big) \\
    P_{k|k} &= (1 - K_k) \cdot P_{k|k-1}
\end{align}
En mode allégé, le Kalman est remplacé par une moyenne mobile exponentielle (\code{LowPassEMA})~:
\begin{equation}
    y[k] = \alpha \cdot y_1[k] + (1 - \alpha) \cdot y[k-1], \quad \text{avec } \alpha = 0{,}25
\end{equation}

\section{Services d'Infrastructure Partagés (\texttt{libs/shared-services/})}

\subsection{Mémoire Pub/Sub par Anneau Circulaire (\texttt{shared::services::memory})}
Les classes \code{MemoryShareComponentServer} et \code{MemoryShareComponentSubscriber} fournissent une communication multi-abonnés sans verrou basée sur un compteur de publication atomique 32 bits et un tampon circulaire en mémoire partagée. L'abonné vérifie si un nouvel échantillon est disponible en comparant son dernier index lu avec l'index de publication exposé par le serveur.

\subsection{Serveur et Client Sockets IPC UNIX (\texttt{shared::services::ipc})}
La classe \code{ServiceIPCSocketServer} crée un socket d'écoute \code{AF\_UNIX} en mode \textbf{non-bloquant} (\code{O\_NONBLOCK}). La méthode \code{process\_pending\_commands()} utilise l'appel système \code{poll()} avec un délai d'attente nul ($0\text{ ms}$) pour consommer les requêtes en instance sans suspendre la boucle de contrôle temps réel critique.

\subsection{Gestionnaire de Configuration (\texttt{shared::services::config})}
La classe \code{ServiceConfig} résout automatiquement le chemin du fichier de configuration en examinant séquentiellement~: \file{/root/<conf>}, \file{./<conf>}, \file{../<conf>}, \file{../../<conf>} et \file{/etc/<conf>}. Elle offre l'écriture synchrone persistante sur disque pour sauvegarder les paramètres calibrés en cours de session.

\chapter{Spécification Exhaustive des Services Logiciels}
\label{chap:services}

\section{\texorpdfstring{\texttt{ServiceManager}}{ServiceManager} (Superviseur Système \& Init PID 1)}

\subsection{Rôle Architectural et Responsabilités}
Le \code{ServiceManager} constitue le processus racine (\code{PID 1}) du système embarqué lorsque le Raspberry Pi démarre sur une image Linux minimale de production, ou un orchestrateur utilisateur en phase de développement. Ses missions sont~:
\begin{enumerate}
    \item Montage initial des pseudo-systèmes de fichiers virtuels indispensables au noyau~: \file{devtmpfs} sur \file{/dev}, \file{proc} sur \file{/proc}, \file{sysfs} sur \file{/sys}, \file{tmpfs} sur \file{/dev/shm} et \file{/run}, \file{devpts} sur \file{/dev/pts}.
    \item Réarmement et nettoyage préventif des ressources IPC résiduelles (\code{unlink("/dev/shm/picar_shm")}, \code{unlink("/tmp/*.sock")}).
    \item Configuration réseau de l'interface de bouclage \code{lo} (127.0.0.1) et attribution d'une adresse IP statique sur le lien Ethernet filaire \code{eth0} (192.168.0.100).
    \item Lancement et surveillance active du serveur SSH léger \code{dropbear} (port 22) pour la maintenance à distance.
    \item Gestion du cycle de vie des cinq microservices opérationnels (\code{FSM}, \code{Sensors}, \code{Motors}, \code{TelemetryService}, \code{PeripheralsService}) avec réapparition automatique en cas de crash.
\end{enumerate}

\subsection{Boucle de Supervision et Paramètres Temporels}
La boucle de surveillance s'exécute à une fréquence de \textbf{10~Hz} (période de scrutation de \textbf{100~ms}) via l'appel système \code{usleep(100000)}~:
\begin{itemize}
    \item \textbf{Collecte asynchrone des processus défunts}~: Appel non bloquant \code{waitpid(-1, \&status, WNOHANG)}.
    \item \textbf{Délai de refroidissement de redémarrage (Cooldown)}~: Un délai minimum de \textbf{100~ms} est imposé entre deux tentatives de relance d'un service défaillant.
    \item \textbf{Protection anti-emballement (Exponential Backoff)}~: Si un service subit plus de 5 échecs consécutifs, le gestionnaire temporise la relance pendant au moins \textbf{2~000~ms} (2 secondes) pour éviter la saturation du CPU par des boucles de crash infinies.
    \item \textbf{Extinction propre (Graceful Shutdown)}~: Lors de la réception de \code{SIGINT} ou \code{SIGTERM}, le superviseur envoie d'abord \code{SIGINT} à tous ses démons enfants, accorde une période de grâce de \textbf{1~500~ms}, puis élimine les processus récalcitrants par \code{SIGKILL} avant de nettoyer la mémoire partagée et les sockets.
\end{itemize}

\subsection{Codes d'Erreur et Conditions de Déclenchement}
\begin{table}[htbp]
\centering
\small
\begin{tabularx}{\textwidth}{lp{4cm}X}
\toprule
\textbf{Condition Détectée} & \textbf{Constat Système} & \textbf{Action \& Code Associé} \\
\midrule
Arrêt inopiné d'un service & \code{WIFEXITED(status) == true} & Log \code{[SUPERVISOR]}, incrémente \code{restart\_count}, re-planifie l'exécution. \\
Mort par signal noyau & \code{WIFSIGNALED(status) == true} & Log du numéro de signal (\code{SIGSEGV}, \code{SIGABRT}), réarmement différé. \\
Instabilité répétée & $\text{restart\_count} \ge 5$ en $< 5\text{ s}$ & Blocage temporaire 2~s, activation du drapeau d'anomalie système. \\
Câble Ethernet débranché & \code{check\_carrier(iface) == false} & Déréférencement de la route IP, log \code{[NET] link DOWN}. \\
Gel d'un processus enfant & Aucun battement depuis $> 500\text{ ms}$ & Déclenchement du drapeau \code{PROCESS\_STALLED} (\reg{0400}). \\
\bottomrule
\end{tabularx}
\caption{Gestion des défaillances par le ServiceManager.}
\label{tab:servicemanager_failures}
\end{table}

% -------------------------------------------------------------
\section{\texorpdfstring{\texttt{FSM}}{FSM} (Machine à États Finis \& Contrôle de Trajectoire)}

\subsection{Rôle Architectural et Composants Internes}
Le service \code{FSM} héberge l'intelligence décisionnelle du robot. Il intègre quatre sous-composants majeurs~:
\begin{enumerate}
    \item \code{FsmNavigator}~: Orchestre le régulateur de suivi de ligne optique (\code{LineFollowerController}), le profileur cinématique trapézoïdal (\code{MotionProfiler}) et la machine d'état de contournement (\code{ObstacleAvoidance}).
    \item \code{FsmShmBridge}~: Assure la lecture non bloquante des capteurs et l'écriture synchrone des consignes moteur/servo dans \file{/dev/shm/picar_shm}.
    \item \code{FsmEventPipe}~: Détecte l'impulsion de départ en provenance du bouton physique via \file{/tmp/peripherals.sock}.
    \item \code{FsmIpcHandler}~: Réceptionne les ordres de réglage dynamique sur \file{/tmp/fsm.sock}.
\end{enumerate}

\subsection{Fréquence d'Exécution et Cadencement Déterministe}
La boucle principale s'exécute à \textbf{50~Hz} (période nominale $T = 20\text{ ms}$). Le respect de la cadence est assuré par un calcul d'échéance absolue anti-dérive (\emph{anti-drift deadline pacing})~:
\begin{lstlisting}[language=C++]
const auto period = std::chrono::nanoseconds(1000000000LL / m_hz);
auto next_deadline = std::chrono::steady_clock::now();
while (m_running.load()) {
    next_deadline += period;
    // ... execution de la logique temps reel ...
    std::this_thread::sleep_until(next_deadline);
}
\end{lstlisting}

\subsection{Paramètres de Configuration (\texttt{fsm.conf})}
\begin{itemize}
    \item \code{HZ} ($50{,}0$)~: Fréquence d'itération de la boucle FSM.
    \item \code{MAX\_ACCELERATION} ($0{,}735\text{ m/s}^2$)~: Accélération maximale admissible (Seuil FastDrink).
    \item \code{STRAIGHT\_CRUISE\_SPEED} ($0{,}35\text{ à } 0{,}45\text{ m/s}$)~: Vitesse cible en ligne droite.
    \item \code{MIN\_CURVE\_SPEED} ($0{,}14\text{ à } 0{,}17\text{ m/s}$)~: Vitesse plancher en courbe de rayon $R = 120\text{ mm}$.
    \item \code{LINE\_KP} ($1{,}1453$) et \code{LINE\_KD} ($0{,}1416$)~: Coefficients proportionnel et dérivé du suivi de ligne.
    \item \code{OBSTACLE\_STOP_DIST\_CM} ($20{,}0\text{ cm}$)~: Distance d'arrêt d'urgence devant un obstacle frontal.
    \item \code{REVERSE\_SPEED} ($0{,}12\text{ m/s}$)~: Vitesse de recul guidé en cas de perte de ruban.
\end{itemize}

\subsection{Codes d'Erreur et Drapeaux Composant (\texttt{shared::fsm::flags})}
\begin{itemize}
    \item \code{DISCOVERABLE} (\reg{0001})~: Service présent et connecté à la SHM.
    \item \code{ANOMALIES\_DETECTED} (\reg{0002})~: Violation de contrainte ou comportement anormal détecté.
    \item \code{WARNING\_DETECTED} (\reg{0004})~: Avertissement mineur (ex: perte momentanée de ligne).
    \item \code{CONFIG\_ERROR} (\reg{0008})~: Fichier de configuration absent ou paramètre hors bornes.
    \item \code{DRIVER\_ERROR} (\reg{0010})~: Défaillance d'accès à la mémoire partagée.
    \item \code{SOCKET\_ERROR} (\reg{0020})~: Échec de création ou d'écoute sur \file{/tmp/fsm.sock}.
    \item \code{SHM\_ERROR} (\reg{0040})~: Segment \file{/dev/shm/picar_shm} introuvable ou verrou bloqué.
    \item \code{STALLED} (\reg{0080})~: Chien de garde capteurs expiré ($> 250\text{ ms}$).
\end{itemize}

% -------------------------------------------------------------
\section{\texorpdfstring{\texttt{Motors}}{Motors} (Actionnement Propulsion \& Direction)}

\subsection{Rôle Architectural et Composants Internes}
Le service \code{Motors} transforme les consignes cinématiques ($v_{\text{target}}, \delta_{\text{target}}$) en signaux matériels physiques (PWM et états logiques GPIO). Il s'appuie sur~:
\begin{enumerate}
    \item \code{MotorActuator}~: Instancie le pilote matériel \code{PiCarHatsMotorDriver} et applique un filtre de limitation de variation angulaire (\emph{slew-rate limiter}) continu pour éviter les à-coups sur le servo de direction ($175^\circ/\text{s}$ en ligne droite, $300^\circ/\text{s}$ en courbe).
    \item \code{MotorShmBridge}~: Lit les consignes de navigation dans la SHM et met à jour la télémétrie locale.
    \item \code{MotorIpcHandler}~: Réceptionne les ajustements de trim de direction sur \file{/tmp/motors.sock}.
\end{enumerate}

\subsection{Fréquence d'Exécution et Cadencement}
La boucle d'actionnement tourne à très haute fréquence (\textbf{100~kHz} par défaut dans \file{motors.conf}), garantissant une latence minimale de réaction et une intégration temporelle continue du slew-rate limiter. La sortie PWM matérielle vers le servo et les moteurs est quant à elle cadencée à \textbf{50~Hz} par le contrôleur PCA9685.

\subsection{Mécanismes de Chien de Garde et Failsafe}
\begin{itemize}
    \item \textbf{Chien de garde FSM (250~ms)}~: Si le compteur \code{heartbeat} de la SHM ne progresse plus pendant 250~ms, le service force immédiatement la consigne de vitesse à $0{,}0\text{ m/s}$ (coupure automatique de sécurité).
    \item \textbf{Arrêt d'urgence immédiat}~: Dès que \code{emergency\_stop} passe à \code{true}, le service invoque \code{emergency\_stop()} qui force les sorties du PCA9685 à 0V via le registre \code{ALL\_LED\_OFF\_H} (\reg{FD}) et bascule les broches de direction TB6612 à l'état BAS.
\end{itemize}

% -------------------------------------------------------------
\section{\texorpdfstring{\texttt{Sensors}}{Sensors} (Acquisition Frontale \& Traitement du Signal)}

\subsection{Rôle Architectural et Composants Internes}
Le service \code{Sensors} effectue la capture de l'environnement physique et son pré-traitement numérique~:
\begin{enumerate}
    \item \code{SensorAcquisition}~: Gère le pilote matériel \code{PiCarHatsSensorDriver} et exécute le pipeline de filtrage ultrasonique (\code{UltrasonicPipeline}).
    \item \code{LineEstimator}~: Calcule l'écart barycentrique continu $e \in [-1{,}0, +1{,}0]$ à partir des noirceurs analogiques des 5 TCRT5000.
    \item \code{SensorShmBridge}~: Écrit les grandeurs filtrées dans la mémoire partagée POSIX et publie la structure \code{SensorsTelemetry} sur l'anneau Pub/Sub.
\end{enumerate}

\subsection{Fréquence d'Échantillonnage et Isolation Acoustique}
Alors que la boucle logicielle scrute la mémoire et le suiveur de ligne jusqu'à \textbf{100~kHz}, \textbf{l'acquisition ultrasonique est strictement isolée à 50~Hz}~:
\begin{lstlisting}[language=C++]
const uint64_t us_interval = (hz > 50) ? static_cast<uint64_t>(hz / 50) : 1;
if (loop_tick % us_interval == 0) {
    float raw = m_driver->read_raw_ultrasonic_cm();
    // Traitement du signal ...
}
\end{lstlisting}
Cette isolation évite de suspendre le thread de perception pendant les 5 à 25~ms nécessaires à la propagation acoustique de l'onde ultrasonore du HC-SR04.

\subsection{Gestion des Anomalies Capteurs}
\begin{itemize}
    \item \textbf{Perte du bus I2C}~: Si plus de 3 lectures I2C consécutives échouent ($\approx 60\text{ ms}$), le capteur de ligne passe en état de repli (toutes les noirceurs à 0) et positionne le drapeau \code{LINE\_LOST}.
    \item \textbf{Timeout acoustique}~: Si aucun écho n'est reçu dans les 5~ms, la valeur $-1{,}0\text{ cm}$ est retournée et \code{consecutive\_us\_errors} est incrémenté.
\end{itemize}

% -------------------------------------------------------------
\section{\texorpdfstring{\texttt{PeripheralsService}}{PeripheralsService} (Interface Matérielle \& Signalisation)}

\subsection{Rôle Architectural et Interfaces}
Le service \code{PeripheralsService} contrôle les interactions physiques locales~:
\begin{itemize}
    \item \textbf{Bouton Start (GPIO 13)}~: Déclenche la transition de l'état \code{Armed} vers le départ de la course.
    \item \textbf{Bouton WiFi (GPIO 6)}~: Déclenche l'activation / bascule de l'interface réseau sans fil.
    \item \textbf{Buzzer (GPIO 5)}~: Fournit des retours auditifs distincts (mélodie d'initialisation, bip de confirmation WiFi, tonalité de départ).
    \item \textbf{Socket IPC UNIX (\file{/tmp/peripherals.sock})}~: Serveur non-bloquant diffusant les notifications d'événements aux services abonnés (notamment \code{FSM} et \code{TelemetryService}).
\end{itemize}

\subsection{Algorithme d'Anti-Rebond Matériel/Logiciel}
La scrutation est cadencée à \textbf{200~Hz} (délai de poll de 5~ms). Un état de bouton n'est validé comme front actif que s'il demeure stable à l'état bas (\emph{Active LOW}) pendant une fenêtre minimale de \textbf{10~ms}~:
\begin{lstlisting}[language=C++]
auto duration = std::chrono::duration_cast<std::chrono::milliseconds>(now - tracker.press_start).count();
if (duration >= DEBOUNCE_MIN_MS && !tracker.triggered) {
    tracker.triggered = true;
    on_trigger(); // Emission de l'evenement
}
\end{lstlisting}

% -------------------------------------------------------------
\section{\texorpdfstring{\texttt{TelemetryService}}{TelemetryService} (Concentrateur \& Télémétrie Temps Réel)}

\subsection{Rôle Architectural et Diffusion Réseau}
Le \code{TelemetryService} assure la surveillance centralisée de l'ensemble du robot~:
\begin{enumerate}
    \item \textbf{Agrégation SHM \& Pub/Sub}~: Échantillonne simultanément la SHM globale et les anneaux Pub/Sub de \code{Sensors}, \code{Motors}, \code{FSM} et \code{Servo}.
    \item \textbf{Calcul Dynamique des Accélérations}~: Dérive en temps réel la vitesse estimée ($a = \frac{\Delta v}{\Delta t}$) et comptabilise le nombre de violations du seuil critique FastDrink ($|a| > 0{,}785\text{ m/s}^2$).
    \item \textbf{Odométrie}~: Intègre la distance parcourue $d = \int |v| \, dt$.
    \item \textbf{Chien de Garde Global}~: Détecte le gel du compteur de heartbeat ($> 500\text{ ms}$) et lève le drapeau \code{PROCESS\_STALLED}.
    \item \textbf{Serveur TCP (Port 91)}~: Permet la connexion d'un banc de diagnostic graphique distant.
    \item \textbf{Émetteur UDP Broadcast (Port 9002)}~: Diffuse les trames de diagnostic en clair à 50~Hz sur le réseau local.
    \item \textbf{Enregistreur USB (\code{UsbLogger})}~: Exporte les métriques de vol sur clé USB externe.
\end{enumerate}

% -------------------------------------------------------------
\section{\texorpdfstring{\texttt{SimBridge}}{SimBridge} (Passerelle Jumeau Numérique Godot 4)}

\subsection{Rôle Architectural et Interchangeabilité}
Le service \code{SimBridge} assure l'interopérabilité complète avec le jumeau numérique sous Godot Engine 4~:
\begin{itemize}
    \item \textbf{Mode SIMULATION}~: Écoute sur le port TCP \textbf{9000}. Le client Godot (\file{network_client.gd}) transmet les données des capteurs virtuels (distances des \code{RayCast3D} et \code{ShapeCast3D}) que le bridge injecte directement dans \file{/dev/shm/picar_shm}. En retour, le bridge transmet à Godot les consignes de vitesse et d'angle de braquage calculées par \code{FSM}.
    \item \textbf{Mode VEHICLE (Réel)}~: Le bridge lit la télémétrie des capteurs physiques et la transmet à Godot uniquement pour la visualisation 3D en temps réel, sans écraser les lectures capteurs matérielles.
\end{itemize}
La synchronisation est verrouillée à \textbf{50~Hz} (20~ms par cycle), correspondant rigoureusement au pas physique (\code{\_physics\_process}) du moteur de jeu.

\chapter{Machine à États Finis (FSM) \& Analyse des Défaillances}
\label{chap:fsm_failures}

\section{Formalisation Mathématique de l'Automate Décisionnel}
L'automate à états finis (\code{FSM}) régit le comportement autonome global du robot selon un sextuplet formel~:
\begin{equation}
    \mathcal{M} = \langle \mathcal{S}, \; \Sigma, \; \mathcal{O}, \; \delta, \; \lambda, \; s_0 \rangle
\end{equation}
où~:
\begin{itemize}
    \item $\mathcal{S}$ représente l'ensemble des 15 états discrets (\code{MachineStates}).
    \item $\Sigma$ est l'ensemble des événements d'entrée (télémétrie capteurs, distances acoustiques, noirceurs optiques, chiens de garde, commandes utilisateur).
    \item $\mathcal{O}$ est le vecteur de consigne d'actionnement en sortie~: $\mathbf{u} = \begin{bmatrix} v_{\text{cible}} & \theta_{\text{cible}} & \text{emergency\_flag} \end{bmatrix}^T$.
    \item $\delta : \mathcal{S} \times \Sigma \to \mathcal{S}$ est la fonction de transition d'état.
    \item $\lambda : \mathcal{S} \times \Sigma \to \mathcal{O}$ est la fonction de sortie cinématique.
    \item $s_0 = \texttt{Init}$ est l'état initial au démarrage.
\end{itemize}

\section{Catalogue Complet des États (\texttt{shared::fsm::MachineStates})}
Le Tableau~\ref{tab:fsm_states} décrit les 15 états définis dans \file{libs/shared-headers/fsm/machine\_states.hpp}.

\begin{table}[htbp]
\centering
\small
\begin{tabularx}{\textwidth}{ccp{4.2cm}X}
\toprule
\textbf{État} & \textbf{Code Binaire} & \textbf{Nom Symbolique} & \textbf{Signification Opérationnelle \& Actionneurs} \\
\midrule
$S_0$ & \reg{00} & \code{Init} & Initialisation des ressources matérielles, IPC et fichiers de configuration ($t < 0{,}1\text{ s}$). $v = 0, \theta = 0$. \\
$S_1$ & \reg{01} & \code{Unarmed} & Calibrage statique des capteurs optiques et repos ($t < 0{,}4\text{ s}$). $v = 0, \theta = 0$. \\
$S_2$ & \reg{02} & \code{Armed} & Système armé prêt pour le départ. Attente de l'impulsion Start ou du délai auto-start. \\
$S_3$ & \reg{03} & \code{LineFollowingForward} & Régulation nominale de trajectoire par PD barycentrique et profil de vitesse adaptatif. \\
$S_4$ & \reg{04} & \code{CollisionRisk} & Décélération progressive d'urgence bornée ($a \le 0{,}735\text{ m/s}^2$) devant un obstacle ($d \le 20\text{ cm}$). \\
$S_5$ & \reg{05} & \code{Stopped} & Véhicule complètement immobilisé ($v < 0{,}03\text{ m/s}$). Analyse de l'obstacle ou fin de course. \\
$S_6$ & \reg{06} & \code{CollisionOccured} & Impact matériel ou obstacle infranchissable. Coupure totale des moteurs. \\
$S_7$ & \reg{07} & \code{RightDeviation} & Manœuvre d'évitement latérale par la droite (+25°) à $0{,}15\text{ m/s}$ (contournement S4). \\
$S_8$ & \reg{08} & \code{LeftDeviation} & Manœuvre d'évitement latérale par la gauche (-25°) à $0{,}15\text{ m/s}$. \\
$S_9$ & \reg{09} & \code{Reverse} & Marche arrière guidée autonome ($v = -0{,}12\text{ m/s}$, recul maximal $< 150\text{ mm}$ - Spécification S5). \\
$S_{10}$ & \reg{0A} & \code{LineSearching} & Contre-braquage directionnel (-24°) pour recouper le ruban noir après évitement. \\
$S_{11}$ & \reg{0B} & \code{LineDetected} & Ruban noir ré-acquis ($D \ge 0{,}55$). Stabilisation de la trajectoire pendant $0{,}2\text{ s}$. \\
$S_{12}$ & \reg{0C} & \code{LineSearchingTimedOut} & Échec de réacquisition de ligne après $4{,}0\text{ s}$. Coupure d'urgence sécuritaire. \\
$S_{13}$ & \reg{0D} & \code{Abort} & Abandon de mission sur demande utilisateur ou opérateur distant. $v = 0, \theta = 0$. \\
$S_{14}$ & \reg{0E} & \code{Failure} & Arrêt d'urgence de sécurité (violation accélération, perte de capteurs, rupture watchdog). \\
\bottomrule
\end{tabularx}
\caption{Définition complète des états de la machine FSM FastDrink.}
\label{tab:fsm_states}
\end{table}

\section{Diagramme de Transition d'États FSM}
La Figure~\ref{fig:fsm_diagram} synthétise le graphe complet des transitions d'états, les conditions de garde temporelles et événementielles associées.

\begin{figure}[htbp]
\centering
\begin{tikzpicture}[
    scale=0.82, every node/.style={transform shape},
    state/.style={circle, draw=blue!70, fill=blue!10, thick, minimum size=1.6cm, align=center, font=\bfseries\footnotesize},
    initstate/.style={state, fill=gray!25, draw=black},
    failstate/.style={state, fill=red!20, draw=red!80},
    trans/.style={-{Stealth[length=2.5mm]}, thick},
    failtrans/.style={-{Stealth[length=2.5mm]}, thick, dashed, red!70!black}
]

% Ligne Supérieure : Initialisation et Armement
\node[initstate] (init) at (-6, 5) {Init\\(0x00)};
\node[state] (unarmed) at (-2, 5) {Unarmed\\(0x01)};
\node[state] (armed) at (2, 5) {Armed\\(0x02)};

\draw[trans] (init) -- node[above, font=\tiny] {$t \ge 0.1\text{ s}$} (unarmed);
\draw[trans] (unarmed) -- node[above, font=\tiny] {$t \ge 0.4\text{ s}$} (armed);

% Ligne Centrale : Suivi Nominal et Arrêt
\node[state, fill=green!15, draw=green!60!black] (line_follow) at (2, 1.5) {LineFollowing\\Forward\\(0x03)};
\node[state] (coll_risk) at (7, 1.5) {CollisionRisk\\(0x04)};
\node[state] (stopped) at (7, -2.5) {Stopped\\(0x05)};

\draw[trans] (armed) -- node[right, font=\tiny] {Start Event} (line_follow);
\draw[trans] (line_follow) -- node[above, font=\tiny] {$d \le 20\text{ cm}$} (coll_risk);
\draw[trans] (coll_risk) -- node[right, font=\tiny] {$v < 0.03\text{ m/s}$} (stopped);

% Fin de parcours en T
\draw[trans] (line_follow) to[out=45, in=135] node[above, font=\tiny] {Barre T $\implies$ Finished} (stopped);
\draw[trans] (stopped) to[out=160, in=-40] node[pos=0.5, below, font=\tiny] {$t \ge 0.5\text{ s}$ \& finished} (armed);

% Ligne Inférieure : Évitement d'obstacle (S4)
\node[state] (dev) at (2, -2.5) {RightDeviation\\(0x07)};
\node[state] (search) at (-2, -2.5) {LineSearching\\(0x0A)};
\node[state] (detected) at (-2, 1.5) {LineDetected\\(0x0B)};

\draw[trans] (stopped) -- node[above, font=\tiny] {$t \ge 0.2\text{ s}$ \& $d \le 25\text{ cm}$} (dev);
\draw[trans] (dev) -- node[above, font=\tiny] {Waypoints ok} (search);
\draw[trans] (search) -- node[right, font=\tiny] {$D_{\max} \ge 0.55$} (detected);
\draw[trans] (detected) -- node[above, font=\tiny] {$t \ge 0.2\text{ s}$} (line_follow);

% Perte de Ligne et Marche Arrière (S5)
\node[state, fill=orange!20, draw=orange!80] (reverse) at (-6, 1.5) {Reverse\\(0x09)};
\draw[trans] (line_follow) to[out=190, in=0] node[below, font=\tiny] {Perte Ligne} (reverse);
\draw[trans] (reverse) to[out=45, in=135] node[above, font=\tiny] {Ligne vue \& $t \ge 0.2\text{ s}$} (line_follow);

% États de Défaillance Terminale
\node[failstate] (timeout_search) at (-2, -5.5) {Searching\\Timeout\\(0x0C)};
\node[failstate] (failure) at (-6, -2.5) {Failure\\(0x0E)};

\draw[failtrans] (search) -- node[right, font=\tiny] {$t \ge 4.0\text{ s}$} (timeout_search);
\draw[failtrans] (reverse) -- node[left, font=\tiny] {$\Delta d \ge 150\text{ mm}$ ou $t \ge 3\text{ s}$} (failure);
\draw[failtrans] (line_follow) to[out=230, in=60] node[pos=0.4, above, font=\tiny] {Watchdog $> 250\text{ ms}$} (failure);
\draw[failtrans] (stopped) to[out=250, in=330] node[pos=0.5, below, font=\tiny] {Urgence SHM} (failure);

\end{tikzpicture}
\caption{Diagramme de transition d'états complet de la machine FSM.}
\label{fig:fsm_diagram}
\end{figure}

\section{Matrice Complète des Défaillances et Déclencheurs de Sécurité}
Le Tableau~\ref{tab:failures_matrix} répertorie l'ensemble des modes de défaillance prévus par le système, les conditions exactes de leur déclenchement ainsi que la procédure d'arrêt d'urgence exécutée.

\begin{small}
\begin{xltabular}{\textwidth}{lp{3.5cm}p{3.5cm}X}
\caption{Matrice d'analyse des modes de défaillance, seuils et mécanismes de repli.} \label{tab:failures_matrix} \\
\toprule
\textbf{Mode de Panne} & \textbf{Condition de Déclenchement} & \textbf{Action Immédiate (Failsafe)} & \textbf{Comportement Système \& Reprise} \\
\midrule
\endfirsthead

\multicolumn{4}{c}{\tablename\ \thetable\ -- \textit{Suite de la page précédente}} \\
\toprule
\textbf{Mode de Panne} & \textbf{Condition de Déclenchement} & \textbf{Action Immédiate (Failsafe)} & \textbf{Comportement Système \& Reprise} \\
\midrule
\endhead

\midrule
\multicolumn{4}{r}{\textit{Suite sur la page suivante...}} \\
\endfoot

\bottomrule
\endlastfoot

\textbf{Perte de Heartbeat Capteurs} & Aucun incrément du compteur \code{heartbeat} de la SHM pendant $\ge 250\text{ ms}$ & FSM bascule immédiatement en \code{Failure} ($S_{14}$). $v = 0, \theta = 0$. & Coupure moteur matérielle via \code{ALL\_LED\_OFF\_H}. Reprise uniquement après réinitialisation complète. \\
\midrule
\textbf{Perte de Heartbeat Navigation} & Le service \code{Motors} constate $> (hz \cdot 0{,}25)$ cycles sans variation de consigne & Moteurs coupent la vitesse ($v_{\text{target}} = 0{,}0$). & Protection contre l'emballement en cas de gel du FSM. Reprise dès rétablissement du flux. \\
\midrule
\textbf{Perte du Bus I2C Ligne} & Échec de lecture continue I2C pendant $\ge 3$ cycles ($\approx 60\text{ ms}$) & Forçage des noirceurs à zéro ($D_i = 0$), bascule en \code{Reverse} ($S_9$). & Le robot tente un recul autonome $< 150\text{ mm}$. Si toujours aucune ligne après 3~s, transition vers \code{Failure}. \\
\midrule
\textbf{Perte de Réflexion Ligne} & Détection de ruban nulle $\sum_{i=0}^4 D_i < 0{,}25$ en suivi nominal & Transition immédiate vers \code{Reverse} ($S_9$). & Vitesse inversée $v = -0{,}12\text{ m/s}$. Réalignement sur le ruban ou arrêt sécurité après $150\text{ mm}$. \\
\midrule
\textbf{Échec Évitement d'Obstacle} & Temps de recherche dans \code{LineSearching} ($S_{10}$) supérieur à $4{,}0\text{ s}$ & Transition vers \code{LineSearchingTimedOut} ($S_{12}$). & Levée de \code{emergency\_stop}, coupure des ponts en H. Intervention manuelle requise. \\
\midrule
\textbf{Gel Processus Système} & Absence de mise à jour SHM détectée par TelemetryService $> 500\text{ ms}$ & Positionnement du bit \code{PROCESS\_STALLED} (\reg{0400}). & Notification sur console de télémétrie et diffusion du code d'anomalie sur le réseau. \\
\midrule
\textbf{Violation Accélération} & Accélération estimée $|a| > 0{,}785\text{ m/s}^2$ ($a > a_{\max} + \epsilon$) & Incrémentation du compteur \code{accel\_violations}, alerte visuelle. & Alerte temps réel. Si rupture physique, signalement immédiat sur le flux télémétrique. \\
\midrule
\textbf{Crash Processus Fils} & Signal noyau \code{SIGSEGV}, \code{SIGBUS} ou \code{SIGABRT} sur un service & Le superviseur \code{ServiceManager} collecte le code via \code{waitpid}. & Pause de refroidissement de 100~ms, puis re-spawn automatique du binaire (backoff 2~s si $\ge 5$ crashs). \\
\midrule
\textbf{Timeout Télémètre Sonar} & Absence de front montant sur GPIO Echo après $5\text{ ms}$ & Retourne la distance invalide $-1{,}0\text{ cm}$. & Rejet de l'échantillon par le filtre médian. Si persistant, signalement de défaillance capteur. \\
\end{xltabular}
\end{small}

\chapter{Contraintes Temporelles, Fréquences \& Analyse Temps Réel}
\label{chap:realtime_timings}

\section{Tableau Récapitulatif des Fréquences et Périodicités Système}
La conception en microservices asynchrones impose un cadencement rigoureux de chaque boucle d'exécution afin de concilier la réactivité face aux événements imprévus et la bande passante limitée des bus matériels. 

Le Tableau~\ref{tab:frequencies_master} synthétise les fréquences nominales, les périodes associées, les mécanismes de synchronisation et les chemins critiques de chaque composant de la pile FastDrink.

\begin{table}[htbp]
\centering
\footnotesize
\begin{tabularx}{\textwidth}{lccccX}
\toprule
\textbf{Composant / Service} & \textbf{Fréquence} & \textbf{Période ($T$)} & \textbf{Gigue Max.} & \textbf{Mécanisme} & \textbf{Chemin Critique / Rôle} \\
\midrule
\textbf{\texttt{ServiceManager}} & 10~Hz & 100~ms & $\pm 5$~ms & \code{usleep()} & Supervision démons, lien Ethernet filaire, Dropbear SSH. \\
\textbf{\texttt{FSM}} (Contrôle) & \textbf{50~Hz} & \textbf{20~ms} & $\pm 0{,}5$~ms & \code{sleep\_until} & Décision, régulateur PD, profileur de vitesse trapézoïdal. \\
\textbf{\texttt{Motors}} (Loop) & 100~kHz & 10~$\mu$s & $\pm 2$~$\mu$s & \code{sleep\_until} & Polling SHM, slew-rate limiter de direction continu. \\
\textbf{PCA9685 PWM} & 50~Hz & 20~ms & Matériel & Hardware & Période PWM servomoteur RC (impulsions 1 à 2~ms). \\
\textbf{\texttt{Sensors}} (Loop) & 100~kHz & 10~$\mu$s & $\pm 2$~$\mu$s & \code{sleep\_until} & Scrutation I2C ligne, calcul barycentrique, mise à jour SHM. \\
\textbf{HC-SR04 Sonar} & \textbf{50~Hz} & \textbf{20~ms} & $\pm 1$~ms & Diviseur $N$ & Impulsion 10~$\mu$s, écoute écho (5 à 25~ms max). \\
\textbf{\texttt{PeripheralsService}}& 200~Hz & 5~ms & $\pm 0{,}5$~ms & \code{poll(5ms)} & Échantillonnage boutons (anti-rebond $\ge 10$~ms), buzzer. \\
\textbf{\texttt{TelemetryService}} & 50~Hz & 20~ms & $\pm 1$~ms & \code{sleep\_until} & Collecte SHM, dérivation d'accélération, UDP broadcast. \\
\textbf{Console Télémétrie} & 5~Hz & 200~ms & $\pm 10$~ms & Modulo loop & Affichage synthétique des métriques sur terminal console. \\
\textbf{\texttt{SimBridge} (Godot)} & 50~Hz & 20~ms & $\pm 0{,}5$~ms & Socket TCP & Synchronisation avec \code{\_physics\_process} de Godot 4. \\
\bottomrule
\end{tabularx}
\caption{Table maîtresse des fréquences d'exécution et contraintes temporelles.}
\label{tab:frequencies_master}
\end{table}

\section{Stratégie d'Isolation Temporelle des Périphériques Lents}
L'un des défis majeurs dans la programmation temps réel sous Linux consiste à éviter qu'une opération d'entrée-sortie lente ne bloque une boucle de contrôle rapide.

Le capteur ultrasonique HC-SR04 en est l'illustration parfaite~: la mesure par écho sonore requiert entre $1\text{ ms}$ (pour un obstacle à 17~cm) et $25\text{ ms}$ (pour la portée maximale). Si cette lecture était exécutée directement à l'intérieur de la boucle rapide de \code{Sensors}, le thread subirait un gel imprévisible, rendant impossible la détection haute fréquence de la ligne optique.

La solution mise en œuvre repose sur une \textbf{isolation de fréquence synchrone par diviseur d'horloge}~:
\begin{equation}
    N = \left\lfloor \frac{f_{\text{loop}}}{f_{\text{ultrason}}} \right\rfloor = \left\lfloor \frac{100\,000}{50} \right\rfloor = 2\,000 \text{ cycles}
\end{equation}
La mesure acoustique n'est déclenchée que lorsque $\text{loop\_tick} \pmod N = 0$. Durant les 1999 autres cycles, le service réutilise la dernière distance filtrée tout en continuant d'échantillonner le ruban noir à pleine vitesse sans aucun retard.

\section{Chronogramme d'Exécution Temps Réel}
La Figure~\ref{fig:timing_diagram} illustre le cadencement relatif des processus sur un intervalle type de 40~ms (deux trames de contrôle à 50~Hz).

\begin{figure}[htbp]
\centering
\begin{tikzpicture}[
    scale=0.78, every node/.style={transform shape},
    timeline/.style={->, thick, >=stealth},
    gridline/.style={help lines, dashed, gray!50},
    slot/.style={draw=black!70, fill=blue!15, rounded corners=2pt, minimum height=0.45cm, font=\scriptsize},
    sonarslot/.style={draw=green!60!black, fill=green!20, rounded corners=2pt, minimum height=0.45cm, font=\scriptsize},
    fsmslot/.style={draw=purple!70!black, fill=purple!20, rounded corners=2pt, minimum height=0.45cm, font=\scriptsize},
    motorslot/.style={draw=orange!70!black, fill=orange!20, rounded corners=2pt, minimum height=0.45cm, font=\scriptsize}
]

% Axes temporels
\draw[timeline] (0, 4) -- (13.5, 4) node[right, font=\footnotesize] {$t$ (ms)};
\draw[timeline] (0, 2.7) -- (13.5, 2.7) node[right, font=\footnotesize] {$t$ (ms)};
\draw[timeline] (0, 1.4) -- (13.5, 1.4) node[right, font=\footnotesize] {$t$ (ms)};
\draw[timeline] (0, 0.1) -- (13.5, 0.1) node[right, font=\footnotesize] {$t$ (ms)};

% Labels de canaux
\node[anchor=east, font=\small\bfseries] at (-0.2, 4.3) {Sensors (I2C Ligne)};
\node[anchor=east, font=\small\bfseries] at (-0.2, 3.0) {Sensors (HC-SR04 Sonar)};
\node[anchor=east, font=\small\bfseries] at (-0.2, 1.7) {FSM (Régulation 50~Hz)};
\node[anchor=east, font=\small\bfseries] at (-0.2, 0.4) {Motors (PCA9685 PWM)};

% Lignes de repère temporel
\foreach \x/\t in {0/0, 3/10, 6/20, 9/30, 12/40} {
    \draw[gridline] (\x, -0.3) -- (\x, 4.8);
    \node[above, font=\tiny\bfseries] at (\x, 4.8) {\t~ms};
}

% Activités Sensors I2C (haute fréquence régulière)
\foreach \x in {0.2, 0.8, 1.4, 2.0, 2.6, 3.2, 3.8, 4.4, 5.0, 5.6, 6.2, 6.8, 7.4, 8.0, 8.6, 9.2, 9.8, 10.4, 11.0, 11.6, 12.2} {
    \draw[fill=blue!30, draw=blue!80] (\x, 3.9) rectangle ++(0.35, 0.4);
}

% Activité HC-SR04 Sonar (rafale 50 Hz avec temps de vol)
\draw[sonarslot, minimum width=1.5cm] (0.8, 2.95) rectangle ++(1.4, 0.45) node[midway] {Echo (7~ms)};
\draw[sonarslot, minimum width=1.3cm] (6.8, 2.95) rectangle ++(1.2, 0.45) node[midway] {Echo (6~ms)};

% Activité FSM (déclenchée à chaque frontière 20 ms)
\draw[fsmslot, minimum width=1.6cm] (2.3, 1.65) rectangle ++(1.8, 0.45) node[midway] {PD + Profiler};
\draw[fsmslot, minimum width=1.6cm] (8.3, 1.65) rectangle ++(1.8, 0.45) node[midway] {PD + Profiler};

% Activité Motors PCA9685 (émission PWM hardware 50 Hz)
\draw[motorslot, minimum width=2.4cm] (0.1, 0.35) rectangle ++(2.0, 0.45) node[midway] {PWM Période 1 (20~ms)};
\draw[motorslot, minimum width=2.4cm] (6.1, 0.35) rectangle ++(2.0, 0.45) node[midway] {PWM Période 2 (20~ms)};

% Flèches de synchronisation IPC
\draw[-{Stealth}, thick, purple!80!black] (2.2, 3.8) to[out=-90, in=90] (2.5, 2.1);
\node[right, font=\tiny, purple!80!black] at (2.4, 2.8) {SHM Write};

\draw[-{Stealth}, thick, orange!80!black] (4.1, 1.65) to[out=-90, in=90] (4.5, 0.8);
\node[right, font=\tiny, orange!80!black] at (4.5, 1.2) {Consigne $\mathbf{u}$};

\end{tikzpicture}
\caption{Chronogramme temps réel illustrant l'isolation de fréquence et l'absence d'inversion de priorité.}
\label{fig:timing_diagram}
\end{figure}

\section{Justification Mathématique des Limites Physiques FastDrink}

\subsection{\texorpdfstring{Démonstration du Seuil Critique d'Accélération ($a_{\max} \le 0{,}735\text{ m/s}^2$)}{Démonstration du Seuil Critique d'Accélération (a <= 0.735 m/s2)}}
Considérons une bille de masse $m$ reposant au fond d'une coupelle sphérique de profondeur utile $\Delta h = 1{,}50\text{ mm} = 0{,}0015\text{ m}$ et de rayon horizontal bord-à-centre $d = 20{,}0\text{ mm} = 0{,}020\text{ m}$.

Le travail mécanique fourni par la force d'inertie lors d'une accélération linéaire $a$ sur la distance $d$ pour franchir la bordure s'écrit~:
\begin{equation}
    W_{\text{inertie}} = F_{\text{inertie}} \cdot d = (m \cdot a) \cdot d
\end{equation}
L'énergie potentielle de pesanteur nécessaire pour faire monter la bille sur la hauteur $\Delta h$ est~:
\begin{equation}
    \Delta E_p = m \cdot g \cdot \Delta h
\end{equation}
À la limite du basculement, par le théorème de l'énergie mécanique ($W_{\text{inertie}} = \Delta E_p$)~:
\begin{equation}
    m \cdot a \cdot d = m \cdot g \cdot \Delta h \implies a = g \cdot \frac{\Delta h}{d}
\end{equation}
En retenant l'accélération gravitationnelle standard $g = 9{,}80665\text{ m/s}^2$~:
\begin{equation}
    a_{\max} = 9{,}80665 \times \frac{0{,}0015}{0{,}020} = 0{,}735498\text{ m/s}^2 \approx \mathbf{0{,}735\text{ m/s}^2}
\end{equation}
Toute accélération instantanée supérieure à cette valeur éjecte irrémédiablement la bille.

\subsection{\texorpdfstring{Quantum Maximal de Vitesse par Pas Temporel ($\Delta t = 20\text{ ms}$)}{Quantum Maximal de Vitesse par Pas Temporel (Delta t = 20 ms)}}
Pour la boucle décisionnelle cadencée à 50~Hz ($\Delta t = 0{,}020\text{ s}$), le saut maximal de vitesse permis par le \code{MotionProfiler} entre deux itérations consécutives est rigoureusement borné par~:
\begin{equation}
    \Delta v_{\text{cycle}} \le a_{\max} \cdot \Delta t = 0{,}735 \times 0{,}020 = 0{,}0147\text{ m/s} \quad (1{,}47\text{ cm/s par cycle})
\end{equation}

\subsection{\texorpdfstring{Vitesse Maximale Admissible en Courbe ($R_{\min} = 120\text{ mm}$)}{Vitesse Maximale Admissible en Courbe (Rmin = 120 mm)}}
Lors du franchissement d'un virage de rayon minimal imposé $R_{\min} = 120\text{ mm} = 0{,}12\text{ m}$ (Spécification S3), l'accélération centripète $a_c = \frac{v^2}{R}$ ne doit pas excéder $a_{\max}$~:
\begin{equation}
    v_{\text{virage}} \le \sqrt{a_{\max} \cdot R_{\min}} = \sqrt{0{,}735 \cdot 0{,}12} = \sqrt{0{,}0882} \approx \mathbf{0{,}297\text{ m/s}} \quad (29{,}7\text{ cm/s})
\end{equation}

\subsection{\texorpdfstring{Profil de Freinage Ciblé et Arrêt de Précision ($\pm 30\text{ mm}$)}{Profil de Freinage Ciblé et Arrêt de Précision (+/- 30 mm)}}
Pour satisfaire la spécification \textbf{S1} (immobilisation devant un obstacle à la distance cible $d_{\text{cible}} = 200\text{ mm}$ avec tolérance stricte $\pm 30\text{ mm}$), la vitesse résiduelle en fonction de la distance mesurée $d$ est régie par la loi de décélération quadratique optimale~:
\begin{equation}
    v_{\text{freinage}}(d) = \sqrt{2 \cdot a_{\max} \cdot (d - d_{\text{cible}})}
\end{equation}
Lorsque $|d - d_{\text{cible}}| \le 2\text{ cm}$, la commande commute sur l'état \code{Stopped} sans choc mécanique, garantissant la parfaite stabilité de la bille.


\end{document}
